\chapter{Mass Kernels and Conditional Vacuum Selection}
\label{ch:ch07_vacuum}

The question of this chapter is which Lorentz-invariant mass terms the chiral block admits and what selects a vacuum. It treats
mass at one constituent, the one-constituent vacua and the trace potential (\cref{lem:tracepot,thm:twophases}), the $272$-dimensional
space $\mathcal M$ of mass kernels with its complex structure (\cref{thm:kernelspace}), the selection of the Standard-Model orbit by an
adopted potential at leading logarithm, matter in the selected vacuum, and the Goldstone bosons, antiparticles, baryon number and
topology of the vacuum orbit. The classification is proved. The selection is conditional on the adopted potential and scheme, global
minimality beyond the local result is a conjecture supported by numerical searches, curvature masses are not pole masses, and the
electroweak stage is missing. The normalisation of couplings and masses is fixed in \cref{app:units}.

% --- source: tex/s14c_mass#0 ((head))


% --- source: tex/s14c_mass#1 (What a rest energy breaks, and the two senses of rotation)
\section{What a rest energy breaks, and the two senses of rotation}
\label{sec:massbreaks}

The arena carries a mass in two readings. In the field-strength reading the rest energy of a field point is a function of the frame
vector, studied in \cref{sec:energy}; the duality rotation $e^{\vartheta G}$ generated by the binding root multiplies the complex
invariant $q$ by $e^{2i\vartheta}$, so for $q\neq0$ only $\vartheta\in\pi\Z$ fixes it (\cref{sec:orth}). In the Dirac reading the mass is
the kernel $T=e_1e_2e_3=\gamma^0J$ of the Grassmann bilinear (\cref{prop:kernels1,cor:fullmass}). This section asks what the second
breaks and compares the rest motion it generates with the binding rotation of the frame. In brief: the Dirac mass breaks no generator
of $\so{3,3}$, and what reduces the six-dimensional symmetry to the interior one is the kinetic term; it breaks the chiral and
Pauli--G\"ursey symmetries of the massless kinetic term down to $\uu1_J$, along three explicit directions; and the two senses of the
Dirac rest rotation are exchanged by charge conjugation, those of the frame's binding rotation by parity. Units are $\hbar=c=1$ unless
a symbol is shown.

% --- source: tex/s14c_mass#2 (The Dirac mass: invariant in six dimensions, breaking the ki)
\subsection{The Dirac mass: invariant in six dimensions, breaking the kinetic algebra}
\label{sec:massdirac}

With $\gamma^0=e_1$ and $\gamma^a=e_{a+3}$ the kinetic kernels are $K^\mu=-\gamma^0\gamma^\mu$ (\cref{prop:kernels1}). A Grassmann bilinear
$\psi^{\!\top}M\psi$ is invariant under $\delta\psi=Y\psi$ exactly when $Y^{\!\top}M+MY=0$, and by \cref{cor:fullmass} the antisymmetric
kernels with this property for all fifteen bivectors $Y=e_ae_b$ are the multiples of $T$. Write $\mathrm{Spin}(3,3)$ for the group
generated by the exponentials of the bivectors.

\begin{proposition}[$\mathrm{Spin}(3,3)$ keeps the mass form, not the kinetic forms]
\label{mb:Tinv}
On the spinor module $\R^8$:
\begin{enumerate}[label=(\roman*),leftmargin=2.2em]
\item For every bivector $Y=e_ae_b$ and every real $s$, $(e^{sY})^{\!\top}T\,e^{sY}=T$.
\item An element $Y\in\so{3,3}$ satisfies $Y^{\!\top}K^\mu+K^\mu Y=\Lambda^\mu{}_\nu K^\nu$ for some real $4\times4$ matrix $\Lambda$ if and only if
$Y\in\lint\oplus\R J$, a space of dimension seven.
\end{enumerate}
\end{proposition}

\begin{proof}
(i) By \cref{lem:blade} the trivector $T$ commutes with a bivector that shares two or no indices with $\{1,2,3\}$ and anticommutes with a
mixed one, and by \cref{thm:sig} the former are the antisymmetric bivectors and the latter the symmetric ones. For antisymmetric $Y$,
$Y^2=-\id$ and $e^{sY}=\cos s\,\id+\sin s\,Y$ is orthogonal and commutes with $T$. For symmetric $Y$, $Y^2=\id$, $e^{sY}=\cosh s\,\id+\sinh s\,Y$ is
symmetric, and $(e^{sY})^{\!\top}T=e^{sY}T=Te^{-sY}$. (ii) The conditions are $256$ linear equations in the fifteen bivector coefficients
of $Y$ and the sixteen entries of $\Lambda$. Their solution space is seven-dimensional and projects injectively onto the bivectors. It
contains $\lint$, under which the $K^\mu$ transform covariantly, and $J$, which is antisymmetric and commutes with every $\gamma^0\gamma^\mu$,
so that $J^{\!\top}K^\mu+K^\mu J=[K^\mu,J]=0$; these span seven dimensions.
\end{proof}

Item (i) is \cref{cor:fullmass} carried from the Lie algebra to the group; item (ii) is fixed by the algebra. Within $\mathrm{Spin}(3,3)$
the Dirac mass breaks nothing. What reduces the six-dimensional symmetry to $\lint\oplus\R J$ is the kinetic term, in which the clock
$\gamma^0$ is paired with the spatial triple. The breaking that a Dirac mass does describe lives in the internal symmetries of the
kinetic term. Let $\mathfrak g_{\rm massless}=\R\,\Omega\oplus\operatorname{span}\{J,e_2\Omega,e_3\Omega\}\cong\uu1_\Omega\oplus\su2$ be the
algebra of \cref{prop:gauge1}, of which invariance of the mass form leaves $\R J$.

\begin{proposition}[What the Dirac mass breaks]
\label{mb:massless}
\begin{enumerate}[label=(\roman*),leftmargin=2.2em]
\item On $\mathfrak g_{\rm massless}$ the map $Y\mapsto Y^{\!\top}T+TY$ has rank three and image $\operatorname{span}\{T\Omega,\,Te_2\Omega,\,Te_3\Omega\}$,
contained in the six-dimensional kernel space of the interior Lorentz algebra (\cref{thm:subkernels}). The kernel $T\Omega$ commutes
with $J$ and is the pseudoscalar mass of \cref{prop:massangle}; $Te_2\Omega$ and $Te_3\Omega$ anticommute with $J$ and are Majorana kernels.
\item For $M=T(m_s+m_p\Omega)$ the field equation $K^\mu\partial_\mu\Psi=M\Psi$ of $\half\Psi^{\!\top}(K^\mu\partial_\mu-M)\Psi$, with
$\partial_a\Psi=Jp_a\Psi$, is $\dot\Psi=-JH\Psi$, where
\begin{equation}
\label{mb:eq:Hmp}
H=\gamma^0\gamma^ap_a+\gamma^0(m_s+m_p\Omega),\qquad H^{\!\top}=H,\qquad [H,J]=0,\qquad H^2=\bigl(|\bp|^2+m_s^2+m_p^2\bigr)\id .
\end{equation}
\item The chiral rotation $e^{\xi\Omega}$ preserves every $K^\mu$ and sends $M$ to $(e^{\xi\Omega})^{\!\top}Me^{\xi\Omega}=Me^{2\xi\Omega}$; the stabiliser
of $M\neq0$ in $U(1)_\Omega$ is $\{\pm\id\}$.
\end{enumerate}
\end{proposition}

\begin{proof}
(i) Every element of $\mathfrak g_{\rm massless}$ is antisymmetric, so the map is $Y\mapsto[T,Y]$. By \eqref{eq:bladecomm} $T$ anticommutes
with $\Omega$, $e_2\Omega$ and $e_3\Omega$ and commutes with $J$, which gives the images $2T\Omega$, $2Te_2\Omega$, $2Te_3\Omega$ and $0$. They are
Lorentz-invariant kernels because $T$ is one and $\mathfrak g_{\rm massless}$ commutes with $\lint$. $T$ and $\Omega$ commute with $J$, while
$e_2$ and $e_3$ anticommute with it. (ii) Varying gives $K^\mu\partial_\mu\Psi=M\Psi$, since $K^\mu$ is symmetric and $M$ antisymmetric. With
$K^0=-\id$, and $J$ commuting with $\gamma^0\gamma^a$, $\gamma^0$ and $\Omega$, $\dot\Psi=(K^aJp_a-M)\Psi=-JH\Psi$. The five matrices
$\gamma^0\gamma^a$, $\gamma^0$ and $\gamma^0\Omega=e_4e_5e_6$ are symmetric, square to $\id$ and anticommute pairwise, which gives $H^2$.
(iii) $\Omega\in\mathfrak g_{\rm massless}$ is antisymmetric and anticommutes with $T$, so $e^{-\xi\Omega}Te^{\xi\Omega}=Te^{2\xi\Omega}$; and
since $M^{\!\top}M=(m_s^2+m_p^2)\id$, a nonzero $M$ is invertible, so $Me^{2\xi\Omega}=M$ exactly when $e^{2\xi\Omega}=\id$, that is for $\xi\in\pi\Z$, where $e^{\xi\Omega}=\pm\id$.
\end{proof}

The count of three in (i) is that of \cref{thm:twophases}(a); the explicit broken directions are fixed by the algebra. Their physical
content, that a Dirac mass breaks the axial $U(1)$ and the Pauli--G\"ursey $SU(2)$ of a massless Dirac field down to fermion number
\cite{Pauli57,Gursey58}, is inherited through \cref{thm:realification}. Items (ii) and (iii) are inherited: they are the realification of
$h=\boldsymbol\alpha\cdot\bp+\beta(m_s+i\gamma_5m_p)$ and of the chiral rotation, here derived from the Grassmann kernels; the commuting
first-order form of \cref{sec:dirac} with the same $H$ has the same flow. The two readings share a shape, visible already in the bilinears of \cref{tab:diracmaxwell}. In both, the mass is the
modulus of an order parameter of charge two under a $U(1)$ generated by a complex structure commuting with a Lorentz algebra, $G$ for
$\lG$ and $\Omega$ for $\lint$, and a nonzero modulus leaves only $\{\pm\id\}$. They differ in kind: the Dirac mass is a coupling, or a value
selected by a potential (\cref{thm:twophases}), while the rest energy of a frame vector is a property of a configuration that no
potential selects, the minimum of the energy density being the zero field \eqref{eq:EBbound}.

% --- source: tex/s14c_mass#3 (Particle and antiparticle: the two senses of one rotation)
\subsection{Particle and antiparticle: the two senses of one rotation}
\label{sec:masssenses}

At $\bp=0$ Dirac matter moves by the rest flow $\Psi(\lambda)=e^{-m\lambda T}\Psi_0$ of \cref{mb:restflow}, which on the eigenspaces $e_1=\pm1$ is $e^{\mp m\lambda J}\Psi_0$, with energy $\pm\half m|\Psi|^2$; it is the realification of $\psi=e^{\mp im\lambda}\psi_0$ on $\beta=\pm1$. The arena adds a name. The rest motion is the one-parameter group of the temporal volume element $T$, the unique
$\mathrm{Spin}(3,3)$-invariant mass form (\cref{cor:fullmass,mb:Tinv}), and the clock $\gamma^0=-TJ$ is the quotient of that rotation by the charge
rotation. Read in this way, the single rotation $e^{-m\lambda T}$ splits into the two senses of the charge rotation according to the
eigenvalue of the clock.

The frame has a rotation with two senses as well. With the split-form Lagrangian $\cL_\kappa=\tfrac m2\eta(\alpha_+X,\alpha_+X)$ of
\eqref{eq:splitL}, every solution is $z(\tau)=c_-e^{-i\kappa\tau}+c_+e^{i\kappa\tau}$ with $c_\pm\in\C^3$ (\cref{prop:orbitsEM}). Below
$\mathsf C$, $\varpi$, $\mathsf T$ name the operations. On spinors they are realised by the operators of \cref{thm:C,thm:PT}; on the frame
by $C_f=-\id$, $P_f=\eta$ ($\bE\mapsto-\bE$) and $T_f\colon X(\tau)\mapsto-\eta X(-\tau)$, which combines the reversal of $\tau$ with the field
map $(\bB,\bE)\mapsto(-\bB,\bE)$ of $C_fP_f$, as $T=CP$ on fields in \cref{thm:achiral}.

\begin{proposition}[Which operation reverses the sense]
\label{mb:senses}
\begin{enumerate}[label=(\roman*),leftmargin=2.2em]
\item Every unit vector $C$ of the surplus plane $U$ satisfies $C^2=\id$, $[C,T]=0$, $\{C,J\}=\{C,e_1\}=0$ and $C^{\!\top}H(0)C=-H(0)$. It maps the
solution $e^{-m\lambda J}\Psi_0$ on $e_1=+1$ to the solution $e^{+m\lambda J}C\Psi_0$ on $e_1=-1$ of the same flow. The parities $\gamma^0$ and
$\gamma^0J$ and the time reversal $T_{\rm rev}=e_2e_4e_5e_6$ of \cref{thm:PT} keep the sign of the rest energy.
\item On the frame $P_fGP_f^{-1}=-G$, and on $z=c_-e^{-i\kappa\tau}+c_+e^{i\kappa\tau}$
\[
P_f\colon(c_-,c_+)\mapsto(\bar c_+,\bar c_-),\qquad C_f\colon(c_-,c_+)\mapsto(-c_-,-c_+),\qquad T_f\colon(c_-,c_+)\mapsto(-\bar c_-,-\bar c_+).
\]
\item The sense of the Dirac rest rotation is reversed by $\mathsf C$, $\mathsf{CP}$, $\mathsf{CT}$ and $\mathsf{CPT}$ and kept by $\varpi$,
$\mathsf T$ and $\mathsf{PT}$. The sense of the frame's binding rotation is reversed by $\varpi$, $\mathsf{CP}$, $\mathsf{PT}$ and $\mathsf{CPT}$
and kept by $\mathsf C$, $\mathsf T$ and $\mathsf{CT}$.
\end{enumerate}
\end{proposition}

\begin{proof}
(i) Write $C=\cos\sigma\,e_2+\sin\sigma\,e_3$. By \eqref{eq:bladecomm} it anticommutes with $e_1$, with $J$ and with the $\gamma^a$, and
commutes with $T$; it is symmetric with $C^2=\id$. Hence $C^{\!\top}e_1C=-e_1$, $C$ maps the eigenspace $e_1=+1$ onto $e_1=-1$, and $C$ commutes
with $e^{-m\lambda T}$. The operators $\gamma^0=e_1$ and $\gamma^0J=T$ commute with $e_1$ and with $T$, so they preserve each eigenspace and the
flow. $T_{\rm rev}$ commutes with $e_1$ and anticommutes with $J$ and $T$, so $T_{\rm rev}\Psi(-\lambda)=e^{-m\lambda T}T_{\rm rev}\Psi_0$ solves the
same flow on the same eigenspace. (ii) $\eta G\eta=G^{\!\top}=-G$. On $z$, $P_f$ acts as $z\mapsto\bar z$, $C_f$ as $z\mapsto-z$, and $T_f$ as
$z(\tau)\mapsto-\overline{z(-\tau)}$; substituting the general solution gives the three maps. (iii) Each operation either exchanges the two
senses or keeps both, so the signs multiply under composition; the single operations are (i), (ii) and \cref{thm:PT}.
\end{proof}

Item (i) is inherited, realised by the arena's charge conjugation and time reversal (\cref{thm:C,thm:PT}). Items (ii) and (iii) are
fixed by the algebra; (ii) makes explicit, on the two senses, the action of the operations of \cref{thm:EBsplit}, by which $P_f$ maps
$\cL_\kappa$ to $\cL_{-\kappa}=\cL_\kappa+2m\kappa\,\tfrac{d}{d\tau}(\bE\cdot\bB)$, a symmetry up to a total derivative, while $C_f$ and $T_f$ leave $\cL_\kappa$ invariant.

For Dirac matter the two senses are particle and antiparticle. They are solutions of one Lagrangian, of angular frequency $m$, each with
energy of definite sign, and they are exchanged by $C$, a symmetry of the Grassmann Lagrangian, since it preserves the mass form and every
$K^\mu$ (\cref{thm:C}); on the commuting first-order form $\cL_1$ it reverses the sign of the action, because $CJC^{-1}=-J$ and
$CH(\bp)C^{-1}=-H(-\bp)$. For the frame the two senses are the kernels of the two factors of $\alpha_-\alpha_+=\partial_\tau^2+\kappa^2$. One
carries no momentum and the other all of it, and each alone carries zero energy, since the energy is $2m\kappa^2\Re(c_-\cdot c_+)$
(\cref{prop:orbitsEM}). They are exchanged by parity, which is a symmetry only up to a total derivative, and their frequency
$\kappa=c/\ell$ \eqref{eq:kappa} is a conversion rate, not a mass. In the interior $m$ and $\kappa$ are independent inputs: $m$ multiplies the
Lagrangian and $\kappa$ sets the period. Reading the frame's $\hbar\kappa$ as $mc^2$ would therefore be an added identification, not a result.
That the frame's binding is not a particle--antiparticle pair in the Dirac sense is our reading of these facts. In both cases the
reversal is $\mathsf{CPT}$-odd.

% --- source: tex/s14c_mass#4 (Two vacua)
\subsection{Two vacua}
\label{sec:masstwovacua}

The Fock vacuum of one Dirac momentum is unique and neutral, with energy $-2E$, and charge conjugation lifts to a unitary operator on
the Fock space that anticommutes with the charge and carries $H(\bp)$ to $H(-\bp)$ (\cref{thm:spectra,prop:Cstates}). At rest its energy,
$-2m$ ($-2mc^2$ in restored units), is minus half the sum of the rest energies of the four modes. Only in that sense does the vacuum
carry the rest energy with the opposite sign, and only this vacuum does.

The frame has a different vacuum. On a coordinate subspace with $k$ temporal and $l$ spatial directions the zero-point energy of the
interior oscillators is $\tfrac12\hbar\kappa(k-l)$, and it vanishes on $\R^{3,3}$ (\cref{thm:subfs}). The full frame has no state of
lowest energy for this to be the energy of.

Quantised mode by mode with positive norm, the energy of the full frame has the spectrum $\hbar\kappa\Z$, and every state has one of lower energy (\cref{mb:noground}).

The zero-point values restate \cref{thm:subfs}, and \cref{mb:noground} makes explicit for the full space what \cref{sec:subfields} records for the mixed planes. The two
vacua differ in kind: the Dirac vacuum has energy $-2E$ per momentum and is bounded below, while the frame's zero-point energy cancels
between positive temporal and negative spatial oscillators, and there is no ground state for it to be the energy of.

% --- source: tex/s14c_mass#5 (Goldstone directions across the tower)
\subsection{Goldstone directions across the tower}
\label{sec:massgoldstone}

\Cref{mb:tab:goldstone} collects the broken directions at the stages where the paper computes them, at $n$ constituents, counted for the specified generators:
at one constituent the four of $\mathfrak g_{\rm massless}$, at three constituents the $21$ Pati--Salam generators and the central $J$ and
$J\Gamma$ (\cref{prop:global}). One element runs through all of them.

\begin{table}[t]
\centering
\small
\begin{tabular}{@{}>{\raggedright\arraybackslash}p{2.4cm}>{\raggedright\arraybackslash}p{2.3cm}>{\raggedright\arraybackslash}p{1.7cm}>{\raggedright\arraybackslash}p{4.4cm}>{\raggedright\arraybackslash}p{3.0cm}@{}}
\toprule
stage & order parameter & kept & broken directions & source\\
\midrule
$n=1$, Dirac & $T=\gamma^0J$ & $\uu1_J$ & $3$, global: $T\Omega$, $Te_2\Omega$, $Te_3\Omega$ & \cref{thm:twophases}(a), \cref{mb:massless}\\
$n=1$, Majorana & kernel on $\Gamma=+1$ & $\R(J-J\Gamma)$ & $3$, global & \cref{thm:twophases}(b)\\
$n=3$, Pati--Salam & $\Phi_0$ ($\nu_R\nu_R$) & $14$ of $23$ & $9$, eaten by $Z'$, $W_R^\pm$ and six $X$ & \cref{thm:nomajoron}\\
$n=3$, electroweak & $\Phi_0$ and bidoublet & $10$ of $23$ & $12$ eaten, $1$ axion-like ($1-\Gamma$) & \cref{thm:axion}\\
\bottomrule
\end{tabular}
\caption{Broken directions along the tower.}
\label{mb:tab:goldstone}
\end{table}

\begin{lemma}[One charge across the scales]
\label{mb:oneminusGamma}
Let $J$ be the complex structure and $\Gamma$ the chirality, at one constituent or at three, so that $[J,\Gamma]=0$, and let $J-J\Gamma$ be the generator of charge $1-\Gamma$.
\begin{enumerate}[label=(\roman*),leftmargin=2.2em]
\item If a kernel $M\neq0$ commutes with $J$ and anticommutes with $\Gamma$ (Dirac type), then $[J-J\Gamma,M]=-2J\Gamma M\neq0$.
\item If $M$ lives on $\Gamma=+1$, that is $\Gamma M=M\Gamma=M$, and anticommutes with $J$, then $[J-J\Gamma,M]=0$.
\item At one constituent $J-J\Gamma=J+\Omega$.
\end{enumerate}
\end{lemma}

\begin{proof}
$J$ commutes with $\Gamma$. (i) $[J,M]=0$ and $MJ\Gamma=-J\Gamma M$, so $[J\Gamma,M]=2J\Gamma M$, which is nonzero because $J\Gamma$ is invertible.
(ii) $[J,M]=2JM$, and $[J\Gamma,M]=JM-MJ\Gamma=JM+JM\Gamma=2JM$. (iii) $\Gamma=J\Omega$, so $J\Gamma=J^2\Omega=-\Omega$.
\end{proof}

The lemma is inherited: a charge that vanishes on one chirality is kept by a Majorana mass of that chirality and broken by every mass
that pairs opposite chiralities. What the arena adds is the identification across scales: one element of the algebra, $J-J\Gamma$, is
kept by a Majorana mass on $\Gamma=+1$ and broken by every Dirac-type mass, both at one constituent and at three. At one constituent it is
$J+\Omega$; the Majorana phase keeps it, and the Dirac mass $T$, which commutes with $J$ and anticommutes with $\Gamma$, breaks it
(\cref{thm:twophases}). At three constituents $\Phi_0$ lives on $\Gamma=+1$, since $\Phi_0^{\!\top}\Phi_0=P_\nu$ and $P_\nu$ projects into the eigenspace $\Gamma=+1$, and anticommutes with $J$
(\cref{thm:kernelspace}); item (ii) is therefore the one-line reason for the equality $[J,\Phi_0]=[J\Gamma,\Phi_0]$ in \cref{thm:nomajoron}. The bidoublet kernels are of Dirac type (\cref{thm:kernelspace}), so every nonzero bidoublet value breaks
$J-J\Gamma$ by item (i); at the generic electroweak vacuum, with both bidoublet values nonzero, its direction escapes the gauge orbit (\cref{thm:axion}).

% --- source: tex/s14c_mass#6 (What a mass breaks)
\subsection{What a mass breaks}
\label{sec:massverdict}

\begin{table}[t]
\centering
\small
\begin{tabular}{@{}>{\raggedright\arraybackslash}p{3.0cm}>{\raggedright\arraybackslash}p{6.6cm}>{\raggedright\arraybackslash}p{4.6cm}@{}}
\toprule
symmetry & broken by a mass? & where\\
\midrule
interior Lorentz $\lint$ & no: $P^2$ and $T$ are invariant; the little group of a momentum is three-dimensional with or without mass & inherited
\cite{WeinbergQFT}; \cref{sec:energy}\\
frame Lorentz $\lG$ & no: $q$ is invariant; the stabiliser of a nonzero field point is two-dimensional with or without $q$ & inherited
\cite{PenroseRindler84}; \cref{sec:energy}\\
$\mathrm{Spin}(3,3)$ & no: $T$ is invariant; the kinetic term keeps seven dimensions & fixed; \cref{cor:fullmass}, \cref{mb:Tinv}\\
duality $U(1)_G$ & yes, to $\{\pm\id\}$ on Lorentz orbits, exactly when $q\neq0$; for a configuration, not a vacuum & inherited; \cref{sec:orth,sec:energy}\\
axial $U(1)_\Omega$ & yes, to $\{\pm\id\}$ & inherited; \cref{prop:massangle,prop:gauge1}, \cref{mb:massless}\\
Pauli--G\"ursey $\su2$ & yes, to $\uu1_J$ & inherited; \cref{prop:gauge1,thm:twophases}\\
$U(1)_{1-\Gamma}$ & by every Dirac-type mass; not by a Majorana mass on $\Gamma=+1$ & inherited; \cref{thm:twophases,thm:axion}, \cref{mb:oneminusGamma}\\
\bottomrule
\end{tabular}
\caption{What a mass breaks in the arena.}
\label{mb:tab:verdict}
\end{table}

\Cref{mb:tab:verdict} collects the answer. The statement that a mass is what breaks a symmetry is true in the arena for duality in the
frame, in the orbit-space sense and for a configuration rather than a vacuum, and for the chiral and Pauli--G\"ursey symmetries of Dirac
matter. It is false for every spacetime symmetry, the six-dimensional one included: a mass changes the type of a Lorentz stabiliser, never its
dimension, and at a field point what a nonzero $q$ removes is the duality generator (\cref{sec:energy}).

Which term counts as the symmetric part, in the sense of the clause on $c^2$ discussed at the end of \cref{sec:energy}, and which as the breaking depends on which symmetry is meant. For $\mathrm{Spin}(3,3)$ the mass form is invariant and the
kinetic term breaks the symmetry (\cref{mb:Tinv}); for the massless kinetic algebra $\uu1_\Omega\oplus\su2$ it is the mass that cuts it to
$\uu1_J$ (\cref{mb:massless}); and in the frame the generator of duality is the binding root, which enters the interior through the shift
$\kappa G$ but is not an infinitesimal isometry of the split form on which the interior Lagrangian is built (\cref{sec:orth,sec:energy}).
In the closure of the force law the breaking of duality is explicit: with Lorentz-class inertia the duality angle of the momentum field
moves in a potential whose amplitude vanishes exactly when that field is null, that is when its rest energy vanishes
(\cref{thm:pendulum,cor:constrate}). This restates \cref{sec:energy}; the breaking is explicit, carried by the inertia and not by a vacuum, and the reading of the angle as a
pseudo-Goldstone boson of duality, with its limits, is taken up in \cref{sec:restclosure}.

\section{Candidates for a mass at one constituent}
\label{sec:w:candidatesforamassatonec}

% --- source: tex/s14b_energy#10 (Three candidates for m in the arena)
\subsubsection{\texorpdfstring{Three candidates for $m$ in the arena}{Three candidates for m in the arena}}
\label{sec:energy-candidates}

The arena has three structures that might be called a mass: the temporal volume element $T=e_1e_2e_3$, which is the Dirac mass form
(\cref{cor:fullmass}); the parameter $m$ and the conserved energy of the interior dynamics of \cref{sec:neutral}; and the rest energy of the
momentum field in the closure of \cref{sec:closure}. We take them in turn.

\begin{proposition}[The mass form acts on the frame as parity]
\label{en:parity}
\begin{enumerate}[label=(\roman*),leftmargin=2.2em]
\item $\Ad_T|_V=\eta$ fixes $\bB$ and reverses $\bE$: on the frame it is parity. It commutes with the $L_i$ and anticommutes with $G$ and the
$K_i$, so it normalises $\lG$ without being $\lG$-invariant.
\item $\eta X$ corresponds to $\bar z$ and $q(z)=h(\bar z,z)$, so $E_6=\tfrac12h(z,z)$ and
$E_0=\tfrac12|h(\bar z,z)|=\tfrac12\max_\vartheta\eta(e^{\vartheta G}X,e^{\vartheta G}X)$.
\item As a Grassmann mass kernel on spinors, $T$ is invariant under all of $\so{3,3}$, $\lG$ included, in the sense of \cref{def:subkernel}, and
breaks nothing; only its sign changes, under conjugation by a spatial generator.
\end{enumerate}
\end{proposition}

\begin{proof}
(i) $\Ad_T|_V=\eta$ is \cref{lem:Gclifford}; $\eta G\eta=-G$ blockwise, the $L_i$ preserve each triple and the $K_i$ exchange them. (ii)
$z(\eta X)=\bB-i\bE$, $h(\bar z,z)=\sum_jz_j^2$, and $\eta(e^{\vartheta G}X,e^{\vartheta G}X)=\Re(e^{2i\vartheta}q)$. (iii) is \cref{cor:fullmass}.
\end{proof}

The identities are inherited or restate the paper; their consequence is an obstruction by type. The identification of $m$ with $T$ fails on frame
vectors, where the adjoint action of $T$ is parity, while as a kernel on spinors $T$ is invariant under $\lG$. The Dirac mass lives on the interior
reading of the six directions, which \cref{prop:emexclusive} excludes on the same six directions, and even there it breaks nothing. What survives
is a reading: the rest energy is the overlap $\tfrac12|h(\bar z,z)|$ of a field with its mirror image, and $\Ad_T$, the arena's splitting into
two triples, is the real structure whose stabiliser in $U(3)$ is $U(1)\cdot O(3)$ (\cref{en:u3}(iv)).

The second candidate is the interior. Its frame field and its momentum field are compared in \cref{en:interior}: the momentum field moves by pure duality and keeps $E_6(\pi)$, $E_0(\pi)$ and $\mathbf S(\pi)$, while the rest energy of the frame field oscillates in general, and the parameter $m$ is an overall factor of the Lagrangian.

As an identification of a rest energy, the interior is obstructed. In it $m$ is an overall factor and an input independent of $\kappa$, the pure
duality motion of either sense carries zero energy, and, with $X$ a length and $\kappa=c/\ell$, the energy
$\mathcal E=mc^2\cdot2\Re(c_-\cdot c_+)/\ell^2$ has the form $mc^2$ times a dimensionless amplitude; but the amplitude is indefinite, vanishes on the
pure duality orbits and is selected by nothing, since $\ell$ is free. That $\kappa G$ lies in the $\eta$-self-adjoint summand of
\cref{thm:EAdichotomy} is fixed; nothing in the interior turns $\kappa^2$ into a rest energy, and reading the rate as a mass, $mc^2=\hbar\kappa$, would
be an added identification, not a result.

The third candidate is the closure.

\begin{remark}[The momentum field in the closure]
\label{en:closure}
The closure of \cref{sec:closure}, whose bearing on the rest energy is the subject of \cref{sec:restclosure}, settles the third candidate by
results fixed there. For every Hamiltonian, $E_6(\pi)=\tfrac12|p|^2$, $E_0(\pi)=\tfrac12|w(p)|$, with $w$ as in \eqref{eq:wdef}, and
$\mathbf S(\pi)$ are Casimirs (\cref{prop:LP,cl:frozen}), while the same functions of the frame field oscillate in the clamped interior
(\cref{en:interior}) and, with a free rotor in the rotating regime, Lorentz-class inertia and $E_0(\pi_0)>0$, grow without bound
(\cref{thm:freedrift,cl:nogen}). For Lorentz-class $N=(\nu_1\eta^{-1}+\nu_2G\eta^{-1})/m$ the potential of the rotor angle, which is the
duality angle of the momentum field, has the amplitude $|\nu_1+i\nu_2|\,E_0(\pi_0)/m$ (\cref{thm:pendulum}); for $(\nu_1,\nu_2)\neq(0,0)$ the
rate is constant and nonzero exactly when $E_0(\pi_0)=0$ and the initial rate is nonzero (\cref{cor:constrate}), and if $E_0(\pi_0)>0$ it is
constant only at the equilibria, where it vanishes (\cref{cl:flat}). For the duality-class $N=\id/m$ the frame energy
$\tfrac12p^{\!\top}Np$ is $E_6(\pi)/m$, the undivided energy of the momentum field, and the rate is conserved for all data
(\cref{thm:inertiareq}(iii)).
\end{remark}

Read through a dictionary that is not a theorem, in which the clamped rate $\kappa=c/\ell$, with $\ell$ free, stands for the conversion and
$E_0(\pi)$, which has the dimension of a squared momentum, stands for a rest energy, these results say the following. The momentum field's rate keeps its symmetry, constancy and invariance under the duality translation of the rotor, when the momentum
field is null; its rest part $E_0(\pi)$ is the amplitude that breaks it; and $E_6(\pi)$ is a Casimir and breaks nothing. But the breaking is explicit,
by an inertia that \cref{thm:inertiareq} forbids from also conserving the rate; what is broken is the constancy of the rotor rate, a property of
duality, not Lorentz invariance and nothing carried by $c$; and with the undivided energy as Hamiltonian, $N=\id/m$, the rate is constant at the
price of the drift of the frame in \cref{thm:drift}.

% --- source: tex/s18_subfields#0 ((head))


% --- source: tex/s18_subfields#1 (The vacuum at one constituent)
\section{The vacuum at one constituent}
\label{sec:subfields}

The word \emph{vacuum} can mean several things, and the arena gives each of them a precise form. It can mean the structure an algebra
selects on its module; the value of an order parameter chosen by a potential; the ground state of a classical field; or the Fock
vacuum of a quantised one. At one constituent all four readings can be computed for every coordinate three-plane of $\R^{3,3}$ and for the
full algebra. The exercise is worth doing for two reasons: it shows which vacuum properties belong to a part of the arena and which
need all of it, and it contains, in miniature, the vacuum that will be selected at three constituents.

% --- source: tex/s18_subfields#2 (The subfields)
\subsection{The subfields}

A three-plane $W\subset\R^{3,3}$ generates a subalgebra $\Cl(W)\subset\Cl(3,3)$. Up to $O(3,3)$ only its signature matters, $(3,0)$,
$(0,3)$, $(2,1)$ or $(1,2)$, but relative to the structures already built (the clock $e_1$, the surplus plane $U$, the interior four-plane
$E$, and the pairing $e_i\leftrightarrow e_{i+3}$) the coordinate planes fall into the classes of \cref{tab:subfields}.

\begin{table}[t]
\centering
\small
\begin{tabular}{@{}llll@{}}
\toprule
name & plane & signature & position\\
\midrule
$T$ & $\langle e_1,e_2,e_3\rangle$ & $(3,0)$ & temporal triple: clock and surplus plane\\
$S$ & $\langle e_4,e_5,e_6\rangle$ & $(0,3)$ & spatial triple, inside the interior four-plane\\
$M_{21a}$ & $\langle e_1,e_2,e_4\rangle$ & $(2,1)$ & clock, one surplus direction, one space; one partner pair\\
$M_{21b}$ & $\langle e_2,e_3,e_4\rangle$ & $(2,1)$ & surplus plane and one space; no partner pair\\
$M_{12a}$ & $\langle e_1,e_4,e_5\rangle$ & $(1,2)$ & a $2{+}1$ slice of the interior spacetime; one partner pair\\
$M_{12b}$ & $\langle e_2,e_4,e_5\rangle$ & $(1,2)$ & a surplus direction and two spaces; one partner pair\\
$M_{12c}$ & $\langle e_1,e_5,e_6\rangle$ & $(1,2)$ & a $2{+}1$ slice of the interior spacetime; no partner pair\\
\bottomrule
\end{tabular}
\caption{Representatives of the coordinate three-planes of $\R^{3,3}$, classified by signature and by their position relative to the
clock, the surplus plane, the interior four-plane and the pairing.}
\label{tab:subfields}
\end{table}

% --- source: tex/s18_subfields#3 (The algebraic vacuum)
\subsection{The algebraic vacuum}

Let $\omega_W$ be the ordered product of the three generators of $W$.

\begin{theorem}[Complex structure or projector]
\label{thm:subdichotomy}
For every coordinate three-plane $W$ of signature $(p,q)$, $\omega_W$ commutes with the generators in $W$ and anticommutes with those outside it, and
\[
\omega_W^2=-(-1)^q\,\id,\qquad \omega_W^{\!\top}=-(-1)^q\,\omega_W .
\]
So $\omega_W$ is a complex structure for $T$ and for the $(1,2)$ planes, and a symmetric involution, hence a rank-four orthogonal projector
$\Pi_W=\tfrac12(\id+\omega_W)$, for $S$ and the $(2,1)$ planes. Three of these volume elements are operators met before:
\[
\omega_T=e_1e_2e_3=\gamma^0J,\qquad \omega_S=\gamma^1\gamma^2\gamma^3,\qquad \omega_T\,\omega_S=\Gamma ,
\]
the Dirac mass form, the time-reversal factor of \cref{thm:PT}, and the chirality, which is also the volume element of the full algebra.
\end{theorem}

\begin{proof}
A generator inside $W$ shares one index with $\omega_W$, so by \eqref{eq:bladecomm} the sign is $(-1)^{3-1}=+1$; one outside shares none, and
the sign is $(-1)^3=-1$. The square and the transpose follow from \eqref{eq:transpose} with $r=3$: $(-1)^{3}\prod\eta_{aa}=-(-1)^q$.
\end{proof}

The restriction to coordinate planes matters for the transpose, not for the square. For any nondegenerate three-plane $W$ with an orthonormal basis,
$\omega_W$ commutes with the vectors of $W$, anticommutes with those of $W^\perp$, and $\omega_W^2=-(-1)^q\id$. The transpose rule, however, refers to the
positive inner product of $S$, which is adapted to the splitting $V=V^+\oplus V^-$: it holds for every plane spanned by vectors of $V^+$ and vectors of $V^-$,
since rotations inside the two triples act on $S$ by orthogonal matrices, and it fails for tilted planes. For the plane
$W_t=\langle\cosh t\,e_4+\sinh t\,e_1,\,e_5,\,e_6\rangle$, of signature $(0,3)$, $\omega_{W_t}=\cosh t\,e_4e_5e_6+\sinh t\,e_1e_5e_6$ is the sum of a symmetric and an
antisymmetric blade, so it is not symmetric for $t\neq0$, although $\omega_{W_t}^2=\id$. This is the mechanism behind the statement below that boost-generated
variations leave the orthogonal projectors.

Under $\so{3,3}$ the volume elements $\omega_T$ and $\omega_S$ have six-dimensional stabilisers that contain all six compact rotations, so no
compact bivector moves them; only the nine boosts do, and by \cref{thm:ghost} boost-generated variations leave the orthogonal projectors.
Each mixed volume element, by contrast, has a two-dimensional compact stabiliser and a four-dimensional compact orbit: the mixed planes are
the only subfields whose vacuum projector or complex structure admits admissible fluctuations. Charge conjugation $C=e_2$ exchanges $\Pi_S$
and $\Pi_\Gamma$ with their complements and fixes the $(2,1)$ projectors, and the full volume element $\Gamma$ commutes with all of $\so{3,3}$:
its projector is the rigid, Lorentz-invariant vacuum $\tfrac12(\id\pm\Gamma)$ that will reappear in \cref{sec:photon}.

% --- source: tex/s18_subfields#4 (The order-parameter vacuum)
\subsection{The order-parameter vacuum}

Let $\so W$ play the role of a Lorentz algebra, its antisymmetric bivectors acting as rotations and its symmetric ones as boosts.

\begin{definition}[Mass kernels of a subfield]
\label{def:subkernel}
The kernel space $\mathcal M_W$ is the space of antisymmetric $M\in\End(S)$ with $[M,\text{rotation}]=0$ and $\{M,\text{boost}\}=0$ for the generators of
$\so W$, and the internal algebra $\mathfrak g_W$ is the space of antisymmetric operators commuting with $\so W$.
\end{definition}

Antisymmetry makes $\psi^{\!\top}M\psi$ a nonzero Grassmann bilinear, and the two conditions make it Lorentz invariant, since a rotation is
antisymmetric and a boost symmetric. Both conditions are diagonal on blades, so both spaces are spanned by blades.

\begin{theorem}[Kernel spaces and invariant kernels]
\label{thm:subkernels}
The kernel spaces, the internal algebras and the invariant kernels are:
\begin{center}
\small\upshape\setlength{\tabcolsep}{4pt}
\begin{tabular}{@{}lccll@{}}
\toprule
Lorentz algebra & $\dim\mathcal M_W$ & $\dim\mathfrak g_W$ & $\mathcal M_W$ under $\mathfrak g_W$ & invariant kernel\\
\midrule
$\so T$, $\so S$ & $10$ & $10$ & adjoint of $\spp2$, irreducible & none\\
each mixed $\so W$ & $10$ & $6$ & $9\oplus1$ & the plane's own $2{+}1$ mass form\\
full $\so{3,3}$ & $1$ & $0$ & --- & $\gamma^0J=e_1e_2e_3$\\
interior $\so{1,3}$ & $6$ & $4$ & irreducible & none\\
\bottomrule
\end{tabular}
\end{center}
The invariant kernel of each mixed plane is $\pm\gamma^0_W\omega_W$, with $\gamma^0_W$ the generator of minority signature. For the interior
algebra $\mathfrak g_I=\langle J,\Omega,e_2\Omega,e_3\Omega\rangle=\uu1_\Omega\oplus\su2$ is the massless gauge algebra of \cref{prop:gauge1}.
\end{theorem}

\noindent Every entry is a finite computation on the $64$ blades of $\Cl(3,3)$: the kernel conditions are diagonal on blades, $\mathfrak g_W$ is the
commutant of $\so W$ in the kernel-preserving algebra, and the decomposition and the invariant kernels follow from the commutant of $\mathfrak g_W$ on
$\mathcal M_W$. Executable check: a6-1 (exact symbolic or integer arithmetic; \cref{app:methods}).

\begin{corollary}[The full algebra has one mass term]
\label{cor:fullmass}
The only Grassmann mass kernel invariant under all of $\so{3,3}$ is $\gamma^0J=\omega_T$, and no antisymmetric operator commutes with
$\so{3,3}$. The order parameter of the full algebra is a single real number $r$ in $r\gamma^0J$; it breaks nothing and has no Goldstone
direction, and its two signs are exchanged by conjugation with any spatial generator.
\end{corollary}

The temporal triple therefore supplies more than a clock: its volume element is the Dirac mass form, the unique Lorentz-invariant mass term
of the whole algebra. To see what a potential selects, take on $\mathcal M_W$ the trace potential
\begin{equation}
\label{eq:tracepot1}
V=-\mu^2N+\alpha N^2+\beta Q,\qquad N=\Tr M^{\!\top}M,\qquad Q=\Tr(M^{\!\top}M)^2 .
\end{equation}
Whether this potential has a minimum at all depends on its coefficients; the next lemma gives the complete answer on the
six-dimensional kernel space $\mathcal M$ of the interior $\so{1,3}$.

\begin{lemma}[The trace potential on $\mathcal M$]
\label{lem:tracepot}
Let $\mathcal M$ be the kernel space of the interior $\so{1,3}$, spanned by the blades $e_1e_2$, $e_1e_3$, $e_1e_2e_3$, $e_2e_4e_5e_6$, $e_3e_4e_5e_6$ and
$e_2e_3e_4e_5e_6$. For $M\in\mathcal M$, $M^{\!\top}M=\tfrac N8\id+D$ with $D$ traceless and $D^2=d^2\id$, so $Q=\tfrac18N^2+8d^2$ and, for $M\neq0$,
\[
q:=Q/N^2\in\bigl[\tfrac18,\tfrac14\bigr],
\]
with $q=\tfrac18$ exactly when $M^{\!\top}M\propto\id$ (for instance on the blades, and on $\gamma^0J$) and $q=\tfrac14$ exactly when $M$ has rank four. In the
coordinates $u=(c_1,c_2,c_6)$, $v=(c_4,c_5,c_3)$ of the six blades in the order listed, $N=8(|u|^2+|v|^2)$ and $d=2|u\times v|$, so $\mathcal M\cong\C^3$ with
$z=u+iv$: $q=\tfrac18$ iff $u\parallel v$, and $q=\tfrac14$ iff $z\cdot z=0$, that is $u\perp v$ and $|u|=|v|$. At fixed $N$ each of these two sets is connected and of
real dimension three. Write
$a(q)=\alpha+\beta q$, so that $V=-\mu^2N+a(q)N^2$.
\begin{enumerate}[label=(\roman*),leftmargin=2.2em]
\item $V$ is bounded below on $\mathcal M$ iff $a(\tfrac18)>0$ and $a(\tfrac14)>0$, or $\min\{a(\tfrac18),a(\tfrac14)\}=0$ and $\mu^2\le0$.
\item If $\mu^2\le0$ and both $a(\tfrac18),a(\tfrac14)>0$, the unique minimum is $M=0$: no symmetry is broken.
\item If $\mu^2>0$ and both $a(\tfrac18),a(\tfrac14)>0$, the minimum on each ray of angular type $q$ is at $N=\mu^2/(2a(q))$ with value $-\mu^4/(4a(q))$, and the
global minima are the kernels with $q=q_*$ minimising $a$: $q_*=\tfrac18$ for $\beta>0$, $q_*=\tfrac14$ for $\beta<0$, at $N_*=\mu^2/(2(\alpha+\beta q_*))$. For
$\beta=0$ every direction of $\mathcal M$ is a minimum direction.
\item Otherwise $V$ is unbounded below, or flat along a ray, and has no isolated minimum. For example, at $\mu^2=1$, $\alpha=-1$, $\beta=1$ one has
$a<0$ on the whole range, and on the ray $M=t\,\gamma^0J$, $V=-8t^2-56t^4$.
\end{enumerate}
\end{lemma}

\begin{proof}
The six blades are the antisymmetric blades that commute with the interior rotations and anticommute with the interior boosts (\cref{def:subkernel});
$M^{\!\top}M=\tfrac N8\id+D$ with $D^2=d^2\id$ is a polynomial identity in the six coefficients. Since $M^{\!\top}M\ge0$ has the eigenvalues $\tfrac N8\pm d$, each on four
dimensions, $|d|\le\tfrac N8$, which gives the range of $q$ and its endpoints. For fixed $q$, $V$ is a quadratic polynomial in $N\ge0$, and (i)--(iv) follow by
minimising it; on $t\,\gamma^0J$, $N=8t^2$ and $Q=8t^4$.
\end{proof}

So the Dirac and Majorana phases below need $\mu^2>0$ and $a>0$ at both ends of the range; outside that region the potential selects
nothing. Inside it, the minimum set at $N_*$ is the three-dimensional set of the lemma, and the kernels named in the theorem are representatives of it: in the
Majorana phase the minima are all null $z$, not only the complex line supported on $\Gamma=+1$, and the other minima are obtained from that line by the
elements $e_2\Omega$, $e_3\Omega$ of $\mathfrak g_I$, which anticommute with $\Gamma$.

\begin{theorem}[Two phases of the interior Lorentz algebra]
\label{thm:twophases}
For the interior $\so{1,3}$, with $\mathcal M$ the irreducible six-dimensional module of $\uu1_\Omega\oplus\su2$, let $\mu^2>0$, $\alpha+\tfrac18\beta>0$ and
$\alpha+\tfrac14\beta>0$, the bounded case of \cref{lem:tracepot}(iii):
\begin{enumerate}[label=(\alph*),leftmargin=2.2em]
\item \emph{Dirac phase} ($\beta>0$). The mass form $\gamma^0J$ is a minimum; its stabiliser in $\mathfrak g_I$ is $\langle J\rangle$, so
$\uu1_\Omega\oplus\su2$ breaks to $\uu1_J$ with three Goldstone directions. Fermion number is conserved and particle and antiparticle are
distinct.
\item \emph{Majorana phase} ($\beta<0$). The kernels supported on the $\Gamma=+1$ half form a complex line; each anticommutes with $J$, has
rank four, $Q/N^2=\tfrac14$ and $M^{\!\top}M\propto\tfrac12(\id+\Gamma)$, and its stabiliser in $\mathfrak g_I$ is exactly $\langle J-J\Gamma\rangle$,
the left-handed fermion number, again with three Goldstone directions.
\end{enumerate}
\end{theorem}

The Majorana phase is the one-constituent seed of the vacuum selected in \cref{sec:selection}. There, too, a concentrated phase of the
trace potential selects a Majorana mass on one right-handed line, and the unbroken global charges include $1-\Gamma$. At one constituent there
is no gauge group and the three Goldstone directions are physical, classically; at three constituents the corresponding directions are gauge
directions and are eaten.

% --- source: tex/s18_subfields#5 (The field-strength vacuum)
\subsection{The field-strength vacuum}

Read through \cref{thm:emfield}, a vector $X$ restricted to $W$ is a field strength of restricted type, and the interior dynamics of
\cref{sec:neutral} restricts to it.

\begin{theorem}[Field strengths on the subfields]
\label{thm:subfs}
For the coordinate three-planes and the full space:
\begin{center}
\footnotesize\upshape
\setlength{\tabcolsep}{4pt}
\begin{tabular}{@{}lccccc@{}}
\toprule
 & Lorentz types & $\bE\cdot\bB\equiv0$ & partner pairs & Lorentz algebra kept & zero-point energy\\
\midrule
$T$ & magnetic only & yes & $0$ & $\so3$ & $+\tfrac32\hbar\kappa$\\
$S$ & electric only & yes & $0$ & $\so3$ & $-\tfrac32\hbar\kappa$\\
$M_{21b}$ & magnetic, electric, null & yes & $0$ & $\so{1,2}$ & $+\tfrac12\hbar\kappa$\\
$M_{12c}$ & magnetic, electric, null & yes & $0$ & $\so{1,2}$ & $-\tfrac12\hbar\kappa$\\
$M_{21a}$ & all four & no & $1$ & none & $+\tfrac12\hbar\kappa$\\
$M_{12a}$, $M_{12b}$ & all four & no & $1$ & none & $-\tfrac12\hbar\kappa$\\
full & all four & no & $3$ & $\so{1,3}$ & $0$\\
\bottomrule
\end{tabular}
\end{center}
On every subfield the cross term of the split-form Lagrangian is a total derivative, the equation of motion is $\ddot X=-\kappa^2X$, and the
energy is $\tfrac m2\bigl(\eta(\dot X,\dot X)+\kappa^2\eta(X,X)\bigr)$. The zero-momentum flow $\dot X=-\kappa GX$ stays in $W$ exactly on
$W\cap G^{-1}W$, of dimension $0$ or $2$ for the subfields and $6$ for the full space.
\end{theorem}

On $T$ the energy is positive definite and $X=0$ is a true ground state: a pure magnetic field oscillates about zero with a positive quantum
zero-point energy. On $S$ the energy is negative definite: $X=0$ is a maximum, dynamically stable, and its quantum oscillators are ghost-like,
of negative energy or, if the energy is made positive, negative norm. On the mixed planes the energy is indefinite and there is no ground
state; the zero-point energy is $\tfrac12\hbar\kappa(p-q)$, and for the full $\R^{3,3}$ it vanishes. Neutral signature cancels the zero-point
energy exactly. The mixed planes without a partner pair are the field strengths of electrodynamics in $2{+}1$ dimensions, $M_{12c}$ carrying one
magnetic and two electric components and $M_{21b}$ its dual, closed under an $\so{1,2}$ and orthogonal by construction.

% --- source: tex/s18_subfields#6 (The Fock vacuum)
\subsection{The Fock vacuum}

A subfield with three anticommuting generators can carry a first-order Hamiltonian, and quantising it with Majorana operators as in
\cref{thm:module} gives a Fock space.

\begin{theorem}[A universal Fock vacuum]
\label{thm:subfock}
Let $A$ be an antisymmetric first-order generator built from three anticommuting generators of a subfield, with $A^2=-E^2$, and let $Q$ be a
commuting charge, antisymmetric with $Q^2=-\id$ and balanced against $A$, $\Tr(AQ)=0$. On the sixteen-dimensional Fock space the joint spectrum of
energy and charge is
\[
(-2E,0)^{1},\ (-E,\pm1)^{2+2},\ (0,0)^4,\ (0,\pm2)^{1+1},\ (E,\pm1)^{2+2},\ (2E,0)^1 :
\]
a unique charge-zero ground state at $-2E$, the zero-point energy of four modes, and two particles and two antiparticles. The charge is the
volume element $\omega_W$ for $T$ and the $(1,2)$ planes; for $S$ and the $(2,1)$ planes it must be borrowed from the complementary plane.
\end{theorem}

\begin{proof}
Commuting antisymmetric operators with $A^2=-E^2\id$ and $Q^2=-\id$ are simultaneously block diagonal on four orthogonal planes, on the $k$-th of which
$A=Ea_kJ_2$ and $Q=s_kJ_2$ with $a_k,s_k\in\{\pm1\}$. Each plane is one fermionic mode, and a Fock state is labelled by occupation signs $\epsilon_k=\pm1$, with energy
$\tfrac E2\sum_k\epsilon_ka_k$ and charge $\tfrac12\sum_k\epsilon_ks_k$. Since $\Tr(AQ)=-2E\sum_ka_ks_k$, the balance condition says that $a_ks_k=+1$ on two planes and
$-1$ on the other two, and counting the sixteen sign choices gives the table; in particular the ground state $\epsilon_k=-a_k$ has charge
$-\tfrac12\sum_ka_ks_k=0$.
\end{proof}

The balance condition cannot be dropped. $A=Ee_4$ and $Q=e_4$ satisfy every other hypothesis, and then the Fock energy and charge are proportional, so the
ground state at $-2E$ carries charge $\pm2$. The charges named in the theorem are balanced: $\Tr(\omega_T\,p\cdot e\,\omega_T)=-\Tr(p\cdot e)=0$ on $T$, and
$\Tr(AJ)=0$ for $A=p\cdot(e_4,e_5,e_6)$ on $S$, because a product of distinct generators is traceless.

The subfields differ in the Hamiltonians they carry. $T$ carries a massless Weyl-type Hamiltonian $\omega_T\,p\cdot e$ with an intrinsic charge;
each $(1,2)$ plane a massive Dirac Hamiltonian in $2{+}1$ dimensions with an intrinsic charge; $S$ a real Majorana Hamiltonian $p\cdot(e_4,e_5,e_6)$
with no intrinsic charge; and the $(2,1)$ planes only a single antisymmetric generator, a mass term without momentum dependence.

\begin{example}[Charge conjugation needs the whole structure]
\label{ex:subC}
Call charge conjugation of a charged subfield any operator anticommuting with its charge and commuting with every term of its Dirac generator,
the condition of \cref{thm:C}. For the interior $3{+}1$ theory with charge $J$ the solutions are exactly the surplus plane $\langle e_2,e_3\rangle$.
For $T$, $M_{12a}$, $M_{12b}$ and $M_{12c}$ the solution spaces are eight-dimensional, and for $S$ with charge $J$ four-dimensional. For $T$ the
space contains $e_4,e_5,e_6$: the charge conjugation of the temporal Weyl theory lives in the spatial triple, as that of the four-dimensional
theory lives in the surplus plane. Only the full four-dimensional Dirac structure fixes $C$ up to the $U(1)_J$ rotation.
\end{example}

\Cref{tab:fourreadings} summarises the four readings for the two triples and the full algebra. The temporal triple supplies the mass form and
an intrinsic charge, the spatial triple the time-reversal factor and, as a field strength, the negative-energy half, and their product is the
chirality. Charge conjugation, the alternative between a Dirac and a Majorana phase, and the cancellation of the zero-point energy all require
the full algebra; none of them is settled inside a subfield.

\begin{table}[t]
\centering
\small
\begin{tabular}{@{}>{\raggedright\arraybackslash}p{2.2cm}>{\raggedright\arraybackslash}p{3.6cm}>{\raggedright\arraybackslash}p{3.6cm}>{\raggedright\arraybackslash}p{5.0cm}@{}}
\toprule
reading & temporal $T$ & spatial $S$ & full $\Cl(3,3)$\\
\midrule
algebraic & $\omega_T=\gamma^0J$: complex structure, rigid under rotations & $\omega_S=\gamma^1\gamma^2\gamma^3$: projector, rigid, swapped by $C$ &
$\Gamma=\omega_T\omega_S$: projector onto the half-spinors, invariant under $\so{3,3}$\\
order parameter & adjoint of $\spp2$; every vacuum breaks it, six Goldstone directions & the same & one invariant kernel $\gamma^0J$; with the interior
$\so{1,3}$, a Dirac phase keeping $J$ or a Majorana phase keeping $J-J\Gamma$\\
field strength & magnetic, positive energy, ground state & electric, negative energy & every type, indefinite energy, zero-point energy cancels\\
Fock & charged Weyl, $C$ in the spatial triple & Majorana, charge only from outside & charged Dirac, $C=\langle e_2,e_3\rangle$ up to $U(1)_J$\\
\bottomrule
\end{tabular}
\caption{The vacuum at one constituent in four readings.}
\label{tab:fourreadings}
\end{table}

% --- source: tex/s19_kernels#0 ((head))


% --- source: tex/s19_kernels#1 (The mass kernels of the chiral block)
\section{The mass kernels of the chiral block}
\label{sec:kernels}

We now return to three constituents and to the block $V_*=\operatorname{im}P_*$ of \cref{sec:pstar}, which carries one Pati--Salam family
$(\mathbf4,\mathbf2,\mathbf1)_{\Gamma=-1}\oplus(\mathbf4,\mathbf1,\mathbf2)_{\Gamma=+1}$ with a right-handed neutrino. A vacuum for this matter is
\emph{dynamically selected} if it solves $\delta\Gamma_{\rm eff}/\delta\Phi=0$ with a second variation that is non-negative apart from the zero modes
of broken symmetries, the equations being those of a stated theory. It is \emph{compatible with matter} if, over it, the fermions carry a
well-defined Dirac operator, the unbroken gauge algebra acts consistently on them, anomalies cancel, and the permitted couplings produce
acceptable masses. This and the next three sections answer both questions inside the arena. Three ingredients are inputs, and we name them
where they enter: the projector $P_*$, which is the field content; the promotion of the space of mass kernels to a dynamical field; and, for
the radiative statements, the form of the tree potential at the matching scale.

% --- source: tex/s19_kernels#2 (Notation on the block)
\subsection{Notation on the block}

The colour algebra $\so6\cong\su4$ is spanned by the fifteen bivectors of the colour directions $C=\{10,11,12,16,17,18\}$; with
$w=(E_7,E_{13},E_{14},E_{15})$ and $\tau_i=-\tfrac12w_jw_k$ for $(i,j,k)$ a cyclic permutation of $(1,2,3)$, as in \cref{sec:pstar},
the two weak algebras are $t^L_i=\tfrac12(\id-\chi_W)\tau_i$ and
$t^R_i=\tfrac12(\id+\chi_W)\tau_i$. These $15+3+3=21$ generators and the two central elements $J$ and $J\Gamma$ commute with $P_*$, with the
Lorentz action and with $J$, and on the block $q_A=E_2E_3$ acts as $J$ and $E_8E_9$ as $J\Gamma$. With $H=-(E_{10}E_{16}+E_{11}E_{17}+E_{12}E_{18})$
the charges are
\[
\hat t=-Jt,\qquad T^L_3=-Jt^L_3,\quad T^R_3=-Jt^R_3,\qquad \BmL=-\tfrac13JH,\qquad Y=T^R_3+\tfrac12(\BmL),
\]
and the colour $\su3$ is the part of $\so6$ commuting with $H$. The charge table is that of \cref{prop:PSledger}. Finally $\gamma^0$
commutes with all $91$ internal bivectors, with $J$ and with the rotations, and anticommutes with $\Gamma$ and the boosts, so it maps $V_*$
onto an orthogonal rank-$64$ \emph{mirror} on which the same internal algebra acts with $\Gamma$ reversed (\cref{rem:PSsettles}).

% --- source: tex/s19_kernels#3 (The kernel space)
\subsection{The kernel space}

\begin{definition}[Mass kernels]
\label{def:kernel}
A Lorentz-invariant Grassmann mass kernel on the block is an $M\in\End(V_*)$ with
\[
M^{\!\top}=-M,\qquad [M,P_*\Sigma^{ab}P_*]=0,\qquad \{M,P_*\Sigma^{0a}P_*\}=0 ,
\]
for the interior rotations $\Sigma^{ab}$ and boosts $\Sigma^{0a}=\tfrac12\gamma^0\gamma^a$. Their real vector space is $\mathcal M$, with inner product
$\langle M,N\rangle=\Tr M^{\!\top}N$.
\end{definition}

The definition is forced. For a Grassmann-valued $\psi\in V_*\otimes\Lambda$ the form $\psi^{\!\top}M\psi$ sees only the antisymmetric part of
$M$, and it is invariant under $\delta\psi=X\psi$ iff $X^{\!\top}M+MX=0$; the rotations are antisymmetric and the boosts symmetric, which gives the
two conditions. A kernel is exactly a Lorentz-invariant mass term of the block.

\begin{theorem}[The kernel space]
\label{thm:kernelspace}
\begin{enumerate}[label=(\roman*),leftmargin=2.2em]
\item $\dim_\R\mathcal M=272$. A basis is given by the compressions $P_*BP_*$ of the blades $B=\gamma^0e_I$ and $B=\gamma^1\gamma^2\gamma^3e_I$, $e_I$ an
internal blade, that are antisymmetric and commute with the three involutions defining $P_*$: there are $2176$ such blades and $272$ distinct,
mutually orthogonal compressions, and they span all solutions of \cref{def:kernel}.
\item $\mathcal M$ is a complex vector space for the complex structure $\mathcal C$, with $\mathcal C(M)=JM$ on the Majorana-type kernels and $\mathcal C(M)=JM\Gamma$ on the Dirac-type kernels
of (iii); $\mathcal C$ commutes with the action of the $23$ generators. (Left multiplication by $J$ alone does not preserve $\mathcal M$: on a Dirac-type kernel $JM$ is
symmetric.) As a complex representation of $\su4\oplus\su2_L\oplus\su2_R$, with the real dimensions of the summands under the braces,
\[
\begin{aligned}
\mathcal M\cong{}&\underbrace{(\mathbf{10},\mathbf3,\mathbf1)}_{60}\oplus\underbrace{(\mathbf{10},\mathbf1,\mathbf3)}_{60}\oplus\underbrace{(\mathbf{15},\mathbf2,\mathbf2)}_{120}\\
&\oplus\underbrace{2(\mathbf6,\mathbf1,\mathbf1)}_{24}\oplus\underbrace{(\mathbf1,\mathbf2,\mathbf2)}_{8},
\end{aligned}
\]
the components being identified by their Casimir labels, $(-36,-2,0)$, $(-36,0,-2)$, $(-32,-\tfrac34,-\tfrac34)$, $(-20,0,0)$ and
$(0,-\tfrac34,-\tfrac34)$ respectively. Replacing $\mathcal C$ by $-\mathcal C$ conjugates every label. As a real representation the two components with $\mathbf{10}$ are irreducible of
complex type, while $(\mathbf{15},\mathbf2,\mathbf2)$, $(\mathbf6,\mathbf1,\mathbf1)$ and $(\mathbf1,\mathbf2,\mathbf2)$ are of real type and occur with multiplicities $2$, $4$ and $2$.
\item Every basis kernel either anticommutes with $J$ and commutes with $\Gamma$ (\emph{Majorana type}, pairing fermions of equal chirality, $144$ of
them, spanning $(\mathbf{10},\mathbf3,\mathbf1)\oplus(\mathbf{10},\mathbf1,\mathbf3)\oplus2(\mathbf6,\mathbf1,\mathbf1)$), or commutes with $J$ and anticommutes with
$\Gamma$ (\emph{Dirac type}, pairing opposite chiralities, $128$ of them, spanning $(\mathbf{15},\mathbf2,\mathbf2)\oplus(\mathbf1,\mathbf2,\mathbf2)$).
\item No nonzero kernel is invariant under the $21$ Pati--Salam generators.
\item The kernels invariant under the twelve Standard-Model generators form the plane spanned by $\Phi_0$ and $J\Phi_0$, where
\begin{equation}
\label{eq:Phi0}
\Phi_0=P_\nu\,P_*\,\gamma^0E_3E_9E_{13}E_{14}E_{16}E_{17}E_{18}\,P_*\,P_\nu ,
\end{equation}
$P_\nu$ being the projector of the block $V_*$ onto $\Gamma=+1$, $T^R_3=+\tfrac12$, $\BmL=-1$, the right-handed neutrino (on all of $S_3$
the joint eigenspace of these three operators has rank eight). Moreover $\Phi_0^{\!\top}\Phi_0=P_\nu$, of rank
four, $J$ anticommutes with $\Phi_0$, and
\[
[t^R_3,\Phi_0]=J\Phi_0,\qquad [\tfrac13H,\Phi_0]=-2J\Phi_0,\qquad [t_Y,\Phi_0]=0,
\]
so $\Phi_0$ carries $T^R_3=1$, $\BmL=-2$ and $Y=0$. It lies in the $(\mathbf{10},\mathbf1,\mathbf3)$ component.
\end{enumerate}
\end{theorem}

\begin{proof}[Proof of the identification in (iii)]
A kernel commuting with $\Gamma$ preserves the two chirality halves, which carry $(\mathbf4,\mathbf2,\mathbf1)$ and $(\mathbf4,\mathbf1,\mathbf2)$; one anticommuting with
$\Gamma$ exchanges them. A bilinear in two $(\mathbf4,\mathbf2,\mathbf1)$ is an $\su2_R$ singlet and one in two $(\mathbf4,\mathbf1,\mathbf2)$ an $\su2_L$ singlet, while a
bilinear between the halves is a doublet of both. So $(\mathbf{15},\mathbf2,\mathbf2)$ and $(\mathbf1,\mathbf2,\mathbf2)$ are of the second type and the others of
the first, and $60+60+24=144$ and $120+8=128$ agree with the counts. The dimension, the components with their Casimir labels and types,
and $\Phi_0$ are checked in v2-1 to v2-3 and z3-8.
\end{proof}

\begin{figure}[t]
\centering
\begin{tikzpicture}[x=0.055cm,y=1cm,font=\footnotesize]
\def\kbar#1#2#3#4#5{\fill[#4] (#1,0) rectangle (#2,0.9); \draw (#1,0) rectangle (#2,0.9); \node at ({(#1+#2)/2},0.45) {#3}; \node[below,font=\footnotesize] at ({(#1+#2)/2},0) {#5};}
\kbar{0}{60}{$(\mathbf{10},\mathbf3,\mathbf1)$}{tcol!20}{$60$}
\kbar{60}{120}{$(\mathbf{10},\mathbf1,\mathbf3)$}{tcol!40}{$60$}
\kbar{120}{144}{\scriptsize$2(\mathbf6,\mathbf1,\mathbf1)$}{tcol!10}{$24$}
\kbar{144}{264}{$(\mathbf{15},\mathbf2,\mathbf2)$}{scol!20}{$120$}
\kbar{264}{272}{}{scol!40}{$8$}
\node[anchor=south,font=\footnotesize] at (268,0.95) {$(\mathbf1,\mathbf2,\mathbf2)$};
\draw[decorate,decoration={brace,amplitude=5pt}] (0,1.0) -- node[above=5pt] {Majorana type: anticommute with $J$, $144$} (144,1.0);
\draw[decorate,decoration={brace,amplitude=5pt,mirror}] (144,-0.45) -- node[below=5pt] {Dirac type: anticommute with $\Gamma$, $128$} (272,-0.45);
\fill[gcol] (112,0.16) circle (2.2pt);
\draw[gcol,thick,->] (112,0.16) -- (112,-0.55) node[below,align=center,font=\footnotesize] {$\Phi_0$: the only\\ Standard-Model line};
\end{tikzpicture}
\caption{The $272$-dimensional space of Lorentz-invariant mass kernels of the chiral block (\cref{thm:kernelspace}), drawn to scale. The
Standard-Model invariant kernels form a single complex line, the right-handed neutrino Majorana kernel $\Phi_0$, inside $(\mathbf{10},\mathbf1,\mathbf3)$.}
\label{fig:kernels}
\end{figure}

The complex dimension of $\mathcal M$ is $136=\dim\mathrm{Sym}^2(\C^{16})$: a Lorentz-invariant mass term of sixteen Weyl fermions is a complex
symmetric $16\times16$ matrix. In the language of $\so{10}$, $\mathrm{Sym}^2(\mathbf{16})=\mathbf{10}\oplus\mathbf{126}$ \cite{Slansky}, and the
decomposition (ii), read with the complex structure $\mathcal C$, is the Pati--Salam branching of $\mathbf{10}_\C\oplus\mathbf{126}$ (or of its conjugate for $-\mathcal C$), of complex
dimension $136=\tfrac12\cdot272$; the $(\mathbf{10},\mathbf1,\mathbf3)\subset\mathbf{126}$ is the field which, in
minimal left--right and Pati--Salam models, gives the right-handed neutrino a Majorana mass \cite{MohapatraSenjanovic80,MohapatraSenjanovic81}.
Here this content has been obtained from \cref{def:kernel}, with only the Pati--Salam algebra acting. The only component containing a
Standard-Model singlet is $(\mathbf{10},\mathbf1,\mathbf3)$: the colour-singlet parts of $(\mathbf{10},\mathbf3,\mathbf1)$, $(\mathbf{15},\mathbf2,\mathbf2)$ and $(\mathbf1,\mathbf2,\mathbf2)$ carry
$Y=\pm1$ or $\pm\tfrac12$, and $(\mathbf6,\mathbf1,\mathbf1)$ has none.

\begin{theorem}[Symmetric criticality]
\label{thm:symcrit}
Let $V:\mathcal M\to\R$ be smooth and invariant under the Pati--Salam group acting by $M\mapsto gMg^{-1}$. Then for every $r\in\R$
\[
\nabla V(r\Phi_0)=\frac1{\langle\Phi_0,\Phi_0\rangle}\,\frac{d}{dr}V(r\Phi_0)\,\Phi_0 ,
\]
so $r\Phi_0$ is a critical point of $V$ if and only if $r$ is a critical point of $r\mapsto V(r\Phi_0)$.
\end{theorem}

\begin{proof}
The action is orthogonal. If $h$ fixes $x=r\Phi_0$, invariance gives $h\nabla V(x)=\nabla V(x)$, so $\nabla V(x)$ is fixed by the Standard-Model
group and lies in $\operatorname{span}\{\Phi_0,J\Phi_0\}$. Since $J\Phi_0=[t^R_3,\Phi_0]$ is tangent to the orbit, $\langle\nabla V(x),J\Phi_0\rangle=0$, and
only the radial component remains.
\end{proof}

This is Michel's theorem on critical orbits \cite{Michel71}, the principle of symmetric criticality \cite{Palais79}. Its content here is that no
invariant coupling on the whole kernel space can generate a tadpole for any other component along $\Phi_0$. It says nothing about the second
variation or about other minima, which are the subject of \cref{sec:selection}.

% --- source: tex/s19_kernels#4 (The kernel field and its action)
\subsection{The kernel field and its action}

\begin{assumption}
\label{ass:field}
The kernels are the values of a scalar field $\Phi(x)\in\mathcal M$ with kinetic term $\tfrac12\langle D_\mu\Phi,D^\mu\Phi\rangle$ and a potential $V(\Phi)$
invariant under the gauge algebra, and a single Yukawa coupling $y$ in \eqref{eq:Saction}. The overall normalisation of the kinetic term is a
convention; its use of one trace form on all of $\mathcal M$, the single coupling $y$ and the form of $V$ are choices, restricted below only by symmetry.
\end{assumption}

With a real Grassmann field $\psi(x)\in V_*$, a gauge field $A_\mu$ with values in an algebra $\mathfrak g$ of antisymmetric operators, the kinetic
kernels $K^\mu=P_*\gamma^0\gamma^\mu P_*$ (symmetric, $K^0=\id$ on $V_*$; minus the compressed kernels of \cref{prop:kernels1}, a sign
convention reconciled after \cref{r1:prop:dirac}) and the module metric $\modG_{ab}=\Tr(t_a^{\!\top}t_b)/512$, take
\begin{equation}
\label{eq:Saction}
S=\int d^4x\Bigl[\tfrac12\psi^{\!\top}K^\mu D_\mu\psi+\tfrac y2\psi^{\!\top}\Phi\psi+\tfrac12\langle D_\mu\Phi,D^\mu\Phi\rangle-V(\Phi)-\tfrac14\modG_{ab}F^a_{\mu\nu}F^{b\mu\nu}\Bigr],
\end{equation}
with $D_\mu\psi=\partial_\mu\psi+A_\mu\psi$ and $D_\mu\Phi=\partial_\mu\Phi+[A_\mu,\Phi]$. Every term is Lorentz invariant by \cref{def:kernel}, and gauge
invariant for any $\mathfrak g$ inside the $23$-dimensional faithful algebra of the block. The module metric is $1$ on the colour bivectors and
$\tfrac18$ on $t^L$ and $t^R$.

\emph{Normalisation dictionary.} As written, \eqref{eq:Saction} has unit gauge coupling for generators that are orthonormal in the module metric;
a coupling $g_A$ is restored by $-\tfrac1{4g_A^2}\modG_{ab}F^a_{\mu\nu}F^{b\mu\nu}$. The generators of module norm $\tfrac18$, such as $t^L_i$, $t^R_i$ and the colour
generator $\tfrac14(E_{10}E_{16}-E_{11}E_{17})$, are those of the conventional normalisation $\Tr T_aT_b=\tfrac12\delta_{ab}$ in the fundamental: their charge
operators have eigenvalues $\pm\tfrac12$ on a doublet. The conventionally normalised coupling is therefore
\[
g=\sqrt8\,g_A ,
\]
and every vector mass from \cref{sec:selection} on is quoted in $g$; \eqref{eq:Saction} as written is the case $g_A=1$, $g=2\sqrt2$. The kernel field is normalised
by $\tfrac12\langle D\Phi,D\Phi\rangle$ with $\langle M,N\rangle=\Tr M^{\!\top}N$, and $v:=|\Phi|/\sqrt2$. Ratios of masses do not depend on these choices.

\begin{proposition}[Fermion number]
\label{prop:fermionnumber}
$J$ commutes with the kinetic kernels and with the $21$ Pati--Salam generators. The transformation $\psi\mapsto e^{\theta J}\psi$,
$\Phi\mapsto e^{\theta J}\Phi e^{-\theta J}$, $A\mapsto A$ leaves every term of \eqref{eq:Saction} except $V$ invariant; on $\mathcal M$ it acts as
$\Phi\mapsto e^{2\theta J}\Phi$ on Majorana-type kernels and trivially on Dirac-type ones. In the complex structure $J$ each of the sixteen Weyl
fermions has charge $1$, so this is the fermion number $F=3B+L$, and $4B=F+(\BmL)$.
\end{proposition}

By \cref{prop:PSanom} the two central directions $J$ and $J\Gamma$ are anomalous on this matter, so the gauge algebra can be any anomaly-free
subalgebra of the $23$, and $\su4\oplus\su2_L\oplus\su2_R$ is the largest one containing Pati--Salam. A gauge field coupled to an anomalous current
makes the quantum theory inconsistent \cite{GrossJackiw72,BouchiatIliopoulosMeyer72}, which is why the central directions carry no gauge
field; they remain global symmetries. In the vacuum $r\Phi_0$ the stabiliser among the $23$ generators is fourteen-dimensional: the twelve
Standard-Model generators together with $J+\tfrac13H$ and $J-J\Gamma$, whose charges are $1+(\BmL)=4B$ and $1-\Gamma$, baryon number and twice
the left-handed fermion number, both global.

\begin{proposition}[The coupled field equations]
\label{prop:coupled3}
The Euler--Lagrange equations of \eqref{eq:Saction} are
\[
K^\mu D_\mu\psi+y\Phi\psi=0,\quad D_\mu D^\mu\Phi+\nabla V(\Phi)-\tfrac y2\Pi_\mathcal M(\psi\psi^{\!\top})=0,\quad \modG_{ab}D_\nu F^{b\nu\mu}=-\bigl(j^\mu_a(\psi)+j^\mu_a(\Phi)\bigr),
\]
with $\Pi_\mathcal M$ the orthogonal projection onto $\mathcal M$. At constant $\Phi=r\Phi_0$, $\psi=0$, $A=0$ the first and third hold identically and the second
reduces to $\tfrac{d}{dr}V(r\Phi_0)=0$ by \cref{thm:symcrit}. The same holds with $V$ replaced by the effective potential, which is again
Pati--Salam invariant; at one loop
\begin{equation}
\label{eq:V1}
V_1=\frac1{64\pi^2}\,\mathrm{STr}\Bigl[\mathcal M^4\Bigl(\ln\frac{\mathcal M^2}{\mu_R^2}-c\Bigr)\Bigr],\qquad
\mathrm{STr}=\textstyle\sum_{\rm scalars}+3\sum_{\rm vectors}-2\sum_{\rm Weyl},
\end{equation}
the fermion term being the determinant of the Dirac operator of the first equation \cite{ColemanWeinberg73,Jackiw74}; at stationary points it
does not depend on the gauge fixing \cite{Nielsen75}.
\end{proposition}

The fermions, classically zero, act on the vacuum only through the effective potential, and the matter content enters through the Yukawa
coupling $y$. That there is one such coupling is part of \cref{ass:field}, not a consequence of calling $\Phi$ one field: $\Phi$ is reducible under the
Pati--Salam algebra, and symmetry alone allows an independent Yukawa coupling and an independent kinetic normalisation for each of its five irreducible
components, with a mixing between the two copies of $(\mathbf6,\mathbf1,\mathbf1)$. The same is true of the common gauge coupling of $\su4$ and $\su2_R$ used in
\cref{sec:selection}. In the vacuum the block Hamiltonian is $H(p)=p_iK^i+iy\,r\,\Phi_0$; since $\{K^i,K^j\}=2\delta^{ij}$,
$\{K^i,\Phi_0\}=0$ and $-\Phi_0^2=P_\nu$,
\[
H(p)^2=|p|^2+y^2r^2P_\nu :
\]
one Majorana fermion, the right-handed neutrino, with four real components and mass $|y|r$, and sixty massless real components.
Executable check: a6-3 for the three identities used here (exact symbolic or integer arithmetic; \cref{app:methods}). The Euler--Lagrange
equations are a direct variation of \eqref{eq:Saction}, and \eqref{eq:V1} is the inherited Coleman--Weinberg formula; neither is the subject of a check.

\Cref{thm:symcrit} settles stationarity for the whole coupled quantum system along $\Phi_0$. Which orbit is the \emph{minimum} is decided by the
couplings, and that is the subject of the next section. A non-negative second variation on the $(\mathbf{10},\mathbf1,\mathbf3)$ directions alone does not
imply it on all of $\mathcal M$: fluctuations outside $(\mathbf{10},\mathbf1,\mathbf3)$ with the Standard-Model quantum numbers of $(\mathbf{10},\mathbf1,\mathbf3)$ fluctuations can
mix with them, and \cref{sec:fullK} exhibits such a mixing.

% --- source: tex/s20_selection#0 ((head))


% --- source: tex/s20_selection#1 (Selecting the Standard-Model orbit)
\section{Selecting the Standard-Model orbit}
\label{sec:selection}

The Standard-Model line $\Phi_0$ lies in the $(\mathbf{10},\mathbf1,\mathbf3)$ component, which is small enough to be analysed completely. We first classify
the renormalizable potential on that component and give sufficient and necessary conditions for $\Phi_0$ to be its global minimum; then we add the gauge
and matter fluctuations;
and finally we return to the whole of $\mathcal M$.

% --- source: tex/s20_selection#2 (The tensor model)
\subsection{The tensor model}

Let $\su4\oplus\su2_R$ act on $\C^8=\C^4\otimes\C^2$ by Hermitian generators normalised by $\Tr T_aT_b=\tfrac12\delta_{ab}$ in the fundamental, and on complex
symmetric $8\times8$ matrices by $S\mapsto USU^{\!\top}$. The subspace
\[
\mathcal T=\mathrm{Sym}^2(\C^4)\otimes\mathrm{Sym}^2(\C^2)\subset\mathrm{Sym}^2(\C^8),\qquad \dim_\C\mathcal T=30,\qquad |S|^2=\Tr S^\dagger S,
\]
is invariant and irreducible, isomorphic to $(\mathbf{10},\mathbf1,\mathbf3)$. With $e_1,\dots,e_4$ and $e_\pm$ the standard bases of $\C^4$ and $\C^2$, $\BmL=\operatorname{diag}(\tfrac13,\tfrac13,\tfrac13,-1)$
and $T^R_3=\operatorname{diag}(\tfrac12,-\tfrac12)$, the vectors $e_i\otimes e_+$, $e_4\otimes e_+$, $e_i\otimes e_-$, $e_4\otimes e_-$ ($i\le3$) carry $Y=\tfrac23,0,-\tfrac13,-1$:
they are $u_R$, $\nu_R$, $d_R$, $e_R$, and a component of $S$ is a pair of them. Put $u_0=e_4\otimes e_+$ and $S_0=u_0u_0^{\!\top}$.

\begin{proposition}[From kernels to tensors]
\label{prop:tensor}
The component $\mathcal M_{(\mathbf{10},\mathbf1,\mathbf3)}$ is $J$-invariant of real dimension $60$, and with $J$ as complex structure it is isomorphic to $\mathcal T$ by a
Pati--Salam-equivariant $\iota$ with $\Tr\Phi^{\!\top}\Phi=4|\iota(\Phi)|^2$, unique up to a phase, with $\iota(\Phi_0)=e^{i\varphi}S_0$. For
$S=\iota(\Phi)$, $\Phi^{\!\top}\Phi$ is supported on the $\Gamma=+1$ half and has the spectrum of $S^\dagger S$, each eigenvalue four times, so
\[
N:=\Tr\Phi^{\!\top}\Phi=4|S|^2,\qquad Q:=\Tr(\Phi^{\!\top}\Phi)^2=4\Tr(S^\dagger S)^2,\qquad \frac Q{N^2}=\frac14\,\frac{\Tr(S^\dagger S)^2}{|S|^4}.
\]
\end{proposition}

\noindent Executable check: a8-1 (numerical; \cref{app:methods}). It recomputes that the component is $J$-invariant of dimension $60$, that its
Pati--Salam commutant is spanned by $\id$ and $J$, so that it is irreducible over $\C$ and $\iota$ is unique up to a complex scale, which the normalisation
fixes up to a phase, that $\Phi_0$ lies in it with $\Phi_0^{\!\top}\Phi_0=P_\nu$, and that $\Phi^{\!\top}\Phi$ is supported on $\Gamma=+1$ with every eigenvalue
four times. The identification of these eigenvalues with those of $S^\dagger S$ uses the explicit isomorphism $\iota$ of the original computation, which the
verification suite of this version does not reproduce.

Let $P_1,\dots,P_5$ be the orthogonal projectors of $\mathrm{Sym}^2\mathcal T$ onto its five irreducible components
\[
(\mathbf{35},\mathbf5)\oplus(\mathbf{35},\mathbf1)\oplus(\mathbf{20}',\mathbf5)\oplus(\mathbf{20}',\mathbf1)\oplus(\mathbf{45},\mathbf3),\qquad 175+35+100+20+135=465,
\]
and put $I_k(S)=\langle S\otimes S,P_k(S\otimes S)\rangle\ge0$, so that $\sum_kI_k=|S|^4$.

\begin{theorem}[The renormalizable potential of $(\mathbf{10},\mathbf1,\mathbf3)$]
\label{thm:potential}
The most general renormalizable invariant potential on $\mathcal T$ is
\begin{equation}
\label{eq:Vgeneral}
V(S)=-\mu^2|S|^2+\sum_{k=1}^5\lambda_kI_k(S)+\kappa\,\mathcal D(S)+\bar\kappa\,\overline{\mathcal D(S)},\qquad \lambda_k\in\R,\ \kappa\in\C,
\end{equation}
where the holomorphic quartic is
\[
\mathcal D(S)=\tfrac1{24}\,\epsilon^{abcd}\epsilon^{efgh}P_{ab,ef}(S)P_{cd,gh}(S),\qquad P_{ab,ef}(S)=\varepsilon_{su}\varepsilon_{tv}S_{as,et}S_{bu,fv},
\]
with colour indices $a,\dots,h$ and $\su2_R$ indices $s,t,u,v$. There is no invariant cubic. Moreover:
\begin{enumerate}[label=(\roman*),leftmargin=2.2em]
\item $\mathcal D(e^{i\theta}S)=e^{4i\theta}\mathcal D(S)$, every $I_k$ is phase invariant, and $\mathcal D$ is not a combination of the $I_k$.
\item $\mathcal D$ depends on $S$ only through the $(\mathbf{20}',\mathbf1)$ component of $S\otimes S$, and $|\mathcal D(S)|\le\tfrac38I_4(S)$, with equality at
$S=\id_4\otimes\id_2$.
\item Under fermion number $S\mapsto e^{2i\theta}S$, so the $I_k$ are invariant and $\mathcal D\mapsto e^{8i\theta}\mathcal D$: \eqref{eq:Vgeneral} is invariant
under fermion number if and only if $\kappa=0$. $\mathcal D$ has $\BmL=0$ and fermion number $8$, hence baryon number $2$.
\end{enumerate}
\end{theorem}

The bound in (ii) follows from Schur's lemma. The $(\mathbf{20}',\mathbf1)$ component carries, up to scale, one invariant complex-bilinear symmetric
form $\beta$; $\mathcal D(S)=c\,\beta(w,w)$ with $w$ the $(\mathbf{20}',\mathbf1)$ component of $S\otimes S$; and $\beta^\dagger\beta$ commutes with the irreducible action,
so all singular values of $\beta$ are equal and $|\mathcal D(S)|\le|c|\,\sigma\,|w|^2=|c|\,\sigma\,I_4(S)$. The constant $|c|\sigma$ and the uniqueness of $\beta$ are
computed numerically, $|c|\sigma=\tfrac38$, and equality holds at $\id_4\otimes\id_2$ (v1-5, v1-6). The invariance and the phase degree of $\mathcal D$ are checked in v1-5,
the vanishing of $\mathcal D$ to second order at $S_0$ in v1-7.

The trivial representation occurs in $\mathcal T\otimes\bar{\mathcal T}$, $\mathrm{Sym}^2\mathcal T$, $\mathrm{Sym}^3\mathcal T$, $\mathrm{Sym}^2\mathcal T\otimes\bar{\mathcal T}$,
$\mathrm{Sym}^4\mathcal T$, $\mathrm{Sym}^3\mathcal T\otimes\bar{\mathcal T}$ and $\mathrm{Sym}^2\mathcal T\otimes\mathrm{Sym}^2\bar{\mathcal T}$ with multiplicities $1,0,0,0,1,0,5$ (z4-4, exact),
which is the content of the theorem. The phase of $\kappa$ can be removed by $S\mapsto e^{i\alpha}S$, so only $|\kappa|$ matters. The contraction of four
$(\mathbf{10},\mathbf1,\mathbf3)$ fields with two $SU(4)$ $\epsilon$ tensors is the baryon-number-violating coupling of Pati--Salam models through which diquark
components mediate neutron--antineutron oscillation \cite{MohapatraMarshak80,ChackoMohapatra99}.

Two lemmas organise the quartics. The first is the recoupling identity $\Tr(S^\dagger S)^2=I_1-\tfrac12I_2-\tfrac12I_3+\tfrac14I_4+0\cdot I_5$, obtained from
Schur--Weyl duality on $(\C^4)^{\otimes4}\otimes(\C^2)^{\otimes4}$. The second characterises the orbit of $S_0$.

\begin{lemma}[The highest-weight orbit]
\label{lem:hw}
For $S\in\mathcal T\setminus0$ the following are equivalent: $S\otimes S$ lies in $(\mathbf{35},\mathbf5)$; $\Tr(S^\dagger S)^2=|S|^4$; $\operatorname{rank}S=1$;
$S=(a\otimes b)(a\otimes b)^{\!\top}$ with $a\in\C^4$, $b\in\C^2$; $S$ lies on the orbit $\mathcal O_{\rm hw}$ of $\C^\times S_0$.
\end{lemma}

This is the Kostant--Lichtenstein description of the highest-weight orbit by quadrics \cite{Lichtenstein82} (checked in v1-3). On $\mathcal O_{\rm hw}$ the colour factor of $S$
has rank one, so $\mathcal D$ vanishes there; more precisely $\mathcal D$ and its first and second derivatives vanish at $S_0$, and
\[
\mathcal D(s_0S_0+\delta)=s_0\Bigl(\tfrac13\epsilon_{ijk}\epsilon_{lmn}U_{il}D_{jm}D_{kn}-\tfrac1{12}\epsilon_{ijk}\epsilon_{lmn}M_{il}M_{jm}D_{kn}\Bigr)+\mathcal D(\delta),
\]
with $U$, $M$, $D$ the symmetric $3\times3$ matrices of the colour sextets $u_Ru_R$, $u_Rd_R$, $d_Rd_R$. In the vacuum the holomorphic quartic contributes
nothing to the masses and produces the cubic couplings $\kappa v\,\epsilon\epsilon\,\Delta_{uu}\Delta_{dd}\Delta_{dd}$ and $\kappa v\,\epsilon\epsilon\,\Delta_{ud}\Delta_{ud}\Delta_{dd}$ of
diquarks of hypercharges $\tfrac43,\tfrac13,-\tfrac23$, each of baryon number two. With the Yukawa couplings of the diquarks to quark pairs, here the
single coupling $y$ of \eqref{eq:Saction}, the vertex closes into a six-quark $\Delta B=2$ operator, the Pati--Salam mechanism of neutron--antineutron
oscillation; the present limits on the oscillation time are $8.6\times10^7$\,s for free neutrons \cite{BaldoCeolin94} and $4.7\times10^8$\,s from
neutrons bound in oxygen \cite{SuperK21}.

% --- source: tex/s20_selection#3 (Where the Standard-Model orbit is the global minimum)
\subsection{Where the Standard-Model orbit is the global minimum}

Number the components $k=1,\dots,5$ as above, put $f_k(S)=I_k(S)/|S|^4$ (so $f_k\ge0$, $\sum f_k=1$), $d(S)=|\mathcal D(S)|/|S|^4$, and
\[
\Delta_k=\lambda_k-\lambda_1\ (k\ge2),\qquad \Lambda(S)=\sum_k\lambda_kf_k(S)=\lambda_1+\sum_{k\ge2}\Delta_kf_k(S),\qquad \Lambda_\kappa(S)=\Lambda(S)-2|\kappa|\,d(S).
\]
For a unit vector $\hat S$, $\min_\varphi V(re^{i\varphi}\hat S)=-\mu^2r^2+\Lambda_\kappa(\hat S)r^4$, because the $I_k$ do not depend on the phase and
$2\Re(\kappa e^{4i\varphi}\mathcal D(\hat S))$ reaches $-2|\kappa||\mathcal D(\hat S)|$. With $\mu^2>0$ the potential is bounded below iff $\Lambda_\kappa>0$ on the unit
sphere, the minimum along a ray is $-\mu^4/(4\Lambda_\kappa)$, and the vacuum is the configuration of smallest $\Lambda_\kappa$.

\begin{theorem}[Selection without the holomorphic quartic]
\label{thm:select0}
Let $\kappa=0$, $\mu^2>0$, $\lambda_1>0$ and $\Delta_2,\dots,\Delta_5>0$. The global minima of \eqref{eq:Vgeneral} are exactly the points of $\mathcal O_{\rm hw}$ with
$|S|^2=\mu^2/(2\lambda_1)$, with $V_{\min}=-\mu^4/(4\lambda_1)$. They form one orbit of real dimension nine, and the stabiliser of each point is conjugate to
$\su3\oplus\uu1_Y$ inside $\su4\oplus\su2_R$ with $Y=T^R_3+\tfrac12(\BmL)$; together with $\su2_L$, which acts trivially on $(\mathbf{10},\mathbf1,\mathbf3)$, this is the
Standard-Model algebra.
\end{theorem}

\begin{proof}
$\Lambda\ge\lambda_1$, with equality iff $f_2=\dots=f_5=0$, that is iff $S\in\mathcal O_{\rm hw}$ by \cref{lem:hw}; the unit-norm part of $\mathcal O_{\rm hw}$ is a single
orbit, of dimension $18-9$; the orbit dimension and the stabiliser are checked in v1-4.
\end{proof}

The unbroken $\uu1$ is an output, computed as the stabiliser: the weights of the vacuum, not a choice of embedding, fix the hypercharge. This is the
vacuum of Pati--Salam and left--right models with a $(\mathbf{10},\mathbf1,\mathbf3)$ field \cite{PatiSalam,MohapatraSenjanovic81}; the theorem says for which couplings
the dynamics chooses it.

\begin{theorem}[Selection with the holomorphic quartic]
\label{thm:selectk}
Let $\mu^2>0$ and $\kappa\in\C$.
\begin{enumerate}[label=(\alph*),leftmargin=2.2em]
\item If $\lambda_1,\lambda_2,\lambda_3,\lambda_5>0$ and $\lambda_4>\tfrac34|\kappa|$, the potential is bounded below; if it is bounded below then $\lambda_1>0$ and
$2\lambda_1+\lambda_2+2\lambda_3+\lambda_4>\tfrac34|\kappa|$.
\item If $\lambda_1>0$, $\Delta_2,\Delta_3,\Delta_5>0$ and $\Delta_4>\tfrac34|\kappa|$, the global minima are exactly the points of $\mathcal O_{\rm hw}$ with
$|S|^2=\mu^2/(2\lambda_1)$.
\item If those points are global minima, then $\Delta_2\ge0$, $\Delta_3\ge0$, $2\Delta_2+2\Delta_3+4\Delta_4+9\Delta_5\ge0$ and
$\Delta_2+2\Delta_3+\Delta_4\ge\tfrac34|\kappa|$.
\item The second variation at the Standard-Model orbit does not depend on $\kappa$.
\end{enumerate}
\end{theorem}

At the benchmark couplings $\lambda=(1,1.7,2.3,3.1,2.9)$ the bounds read: $\tfrac43\Delta_4=2.8$, sufficient for selection; $\tfrac43(\Delta_2+2\Delta_3+\Delta_4)=7.2$,
necessary for selection and equal to the computed threshold; $\tfrac43\lambda_4=4.13$, sufficient for boundedness; and $\tfrac43(2\lambda_1+\lambda_2+2\lambda_3+\lambda_4)=15.2$,
necessary for boundedness. Minimisations from random starts are consistent with the transition at $|\kappa|=7.2$: at $|\kappa|=6,7,7.5,8$ the runs end on the
Standard-Model orbit, which is always a local minimum, or on the orbit of $\id_4\otimes\id_2$, which is the higher of the two for $|\kappa|\le7$ and the lower for
$|\kappa|\ge7.5$, and no run ends below both (v3-1; sixteen archived starts per value, numerical): at the benchmark,
$f(\id_4\otimes\id_2)=(\tfrac13,\tfrac16,\tfrac13,\tfrac16,0)$ and $d(\id_4\otimes\id_2)=\tfrac1{16}$, so
$\Lambda_\kappa(\id_4\otimes\id_2)-\Lambda_\kappa(S_0)=0.9-\tfrac18|\kappa|$, which changes sign at $|\kappa|=7.2$.

\begin{theorem}[Second variation at the Standard-Model orbit]
\label{thm:secondvar}
At $S_v=vS_0$ the sixty real fluctuations split into $SU(3)\times U(1)_Y$ multiplets, each with one mass:
\begin{center}
\small
\begin{tabular}{@{}llccl@{}}
\toprule
pair & $SU(3)_Y$ & $\dim_\R$ & $m^2/v^2$ & role\\
\midrule
$\nu_R\nu_R$ (modulus) & $1_0$ & $1$ & $4\lambda_1$ & radial\\
$\nu_R\nu_R$ (phase), $\nu_Re_R$, $u_R\nu_R$ & $1_0$, $1_{-1}$, $3_{2/3}$ & $1+2+6$ & $0$ & eaten\\
$e_Re_R$ & $1_{-2}$ & $2$ & $\tfrac43\Delta_2$ & $\Delta_{ee}$\\
$u_Re_R+d_R\nu_R$ & $3_{-1/3}$ & $6$ & $\Delta_5$ & $\Delta_{ue}$\\
$d_Re_R$ & $3_{-4/3}$ & $6$ & $\tfrac23\Delta_2+\Delta_5$ & $\Delta_{de}$\\
$u_Ru_R$ & $6_{4/3}$ & $12$ & $\tfrac43\Delta_3$ & $\Delta_{uu}$\\
$u_Rd_R$ & $6_{1/3}$ & $12$ & $\tfrac23\Delta_3+\Delta_5$ & $\Delta_{ud}$\\
$d_Rd_R$ & $6_{-2/3}$ & $12$ & $\tfrac29\Delta_2+\tfrac29\Delta_3+\tfrac49\Delta_4+\Delta_5$ & $\Delta_{dd}$\\
\bottomrule
\end{tabular}
\end{center}
Consequently $S_v$ is a non-degenerate local minimum on $\mathcal T$ iff $\lambda_1,\Delta_2,\Delta_3,\Delta_5>0$ and $2\Delta_2+2\Delta_3+4\Delta_4+9\Delta_5>0$, a region that
strictly contains the cone of \cref{thm:select0}; and, independently of the couplings,
\begin{equation}
\label{eq:sumrules}
m_{de}^2=m_{ue}^2+\tfrac12m_{ee}^2,\qquad m_{ud}^2=m_{ue}^2+\tfrac12m_{uu}^2 .
\end{equation}
\end{theorem}

The table is checked against the full $60\times60$ mass matrix at the benchmark and at random couplings (v1-9 to v1-11). The $d_Rd_R$ entry is
$\tfrac19(2\Delta_2+2\Delta_3+4\Delta_4+9\Delta_5)$, whose positivity is the last local-minimum condition stated in the theorem, and with the one-loop gaps of
\cref{thm:oneloop1013} it gives $m_{dd}^2=120-4x$, as in \cref{tab:271}.

\begin{example}[Checking a sum rule]
\label{ex:sumrule}
At the benchmark $\lambda=(1,1.7,2.3,3.1,2.9)$ one has $\Delta=(0.7,1.3,2.1,1.9)$, so in units of $v^2$: $m_{ee}^2=\tfrac43\cdot0.7=0.9\overline3$,
$m_{ue}^2=1.9$, $m_{de}^2=\tfrac23\cdot0.7+1.9=2.3\overline6$, and indeed $1.9+\tfrac12\cdot0.9\overline3=2.3\overline6$. Likewise $m_{uu}^2=\tfrac43\cdot1.3=1.7\overline3$
and $m_{ud}^2=\tfrac23\cdot1.3+1.9=2.7\overline6=1.9+\tfrac12\cdot1.7\overline3$. The full $60\times60$ mass matrix at these couplings and $\mu^2=2$ has the
predicted spectrum to $4\times10^{-15}$.
\end{example}

The colour-sextet components $\Delta_{uu},\Delta_{ud},\Delta_{dd}$ are the diquarks just met, and $\Delta_{ee}$ is the doubly charged scalar of left--right models.
The sum rules \eqref{eq:sumrules} hold for the most general renormalizable potential of a single $(\mathbf{10},\mathbf1,\mathbf3)$, with or without $\kappa$. When one of the
bounds $\Delta_2,\Delta_3,\Delta_5>0$ fails, gradient flows end on configurations of lower energy that break electric charge or colour: the configuration
$(e_4e_4)(e_+e_-)$, for instance, keeps colour, but its unbroken abelian factor is $T^R_3$ instead of $Y$.

% --- source: tex/s20_selection#4 (The trace potential)
\subsection{The trace potential}

The kernel field carries two quartic invariants that do not refer to the Pati--Salam decomposition at all, $N^2$ and $Q=\Tr(\Phi^{\!\top}\Phi)^2$; they are
invariant under every orthogonal change of basis of the block that preserves $\mathcal M$, in particular under fermion number, so the trace potential has
$\kappa=0$.

\begin{theorem}[The trace potential on $(\mathbf{10},\mathbf1,\mathbf3)$]
\label{thm:tracepot}
Let $V_{\rm tr}=-\mu^2N+\alpha N^2+\beta Q$ with $\mu^2>0$. In the variables of \cref{thm:potential}, $\lambda_k=16\alpha+4\beta a_k$ with $a=(1,-\tfrac12,-\tfrac12,\tfrac14,0)$,
so $(\Delta_2,\Delta_3,\Delta_4,\Delta_5)=-4\beta\,(\tfrac32,\tfrac32,\tfrac34,1)$. If $\beta<0$ and $\alpha+\tfrac14\beta>0$, the global minimum is the Standard-Model orbit; if
$\beta>0$ and $\alpha+\tfrac1{32}\beta>0$, the minima are the $S$ with $S^\dagger S=\tfrac18|S|^2\id_8$, and every one of them breaks colour; if $\beta=0$ the potential is
$O(60)$-symmetric and the orbit is chosen by fluctuations.
\end{theorem}

$Q/N^2$ measures how concentrated the Majorana mass matrix is: it equals $\tfrac14$ when one fermion carries all of it and $\tfrac1{32}$ when the eight
right-handed fermions share it equally. A negative $\beta$ rewards concentration on a single line, and the only line on which the dynamics can concentrate
without breaking colour is $\nu_R$. The Standard-Model vacuum is the concentrated phase of the Majorana mass, which is the three-constituent form of
the Majorana phase of \cref{thm:twophases}.

% --- source: tex/s20_selection#5 (The fluctuations decide: back-reaction on (10,1,3)(10,1,3))
\subsection{The fluctuations decide: back-reaction on \texorpdfstring{$(\mathbf{10},\mathbf1,\mathbf3)$}{(10,1,3)}}

Take the tree potential degenerate on $(\mathbf{10},\mathbf1,\mathbf3)$, $\lambda_1=\dots=\lambda_5$ and $\kappa=0$, the case $\beta=0$. Every unit direction of $\mathcal T$ is then a
tree-level minimum, and the orbit is chosen by the one-loop potential \eqref{eq:V1}. Gauge and Yukawa interactions respect fermion number, so they do not
generate $\kappa$ in perturbation theory. Write $L=\ln(\Lambda_{\rm UV}^2/v^2)$. The orbit-dependent part of the leading logarithm is
\[
\delta V=\frac L{64\pi^2}\Bigl[2\sum_{\rm fam}\Tr(M_F^\dagger M_F)^2-3\Tr M_V^4-\Tr M_S^4\Bigr],
\]
where for the degenerate tree potential $\Tr M_S^4$ depends on $|S|$ only. Here and in \cref{sec:fullK} the field is normalised canonically. All radiative statements
of this section share four limitations: they keep the leading logarithm only, with orbit-independent and non-logarithmic terms dropped and their effect not
estimated; vectors are counted with three polarisations, as in Landau gauge, and the scheme dependence away from stationary points is not examined; the
masses are curvature masses of the effective potential, not pole masses; and the couplings are taken equal where stated, as in \cref{ass:field}. By
\cref{prop:tensor}, $\Tr\Phi^{\!\top}\Phi=4|\iota(\Phi)|^2$, so the kinetic term of \cref{ass:field} is $\Tr(DS)^\dagger(DS)$ for $S:=\sqrt2\,\iota(\Phi)$, and from
here on $S$ denotes this field and $v=|S|=|\Phi|/\sqrt2=\sqrt2\,r$, the $v$ of \cref{sec:fullK,sec:matter}. In these variables the right-handed neutrino mass
$|y|r$ of \cref{sec:kernels} reads $(|y|/\sqrt2)\,v$, so the $y$ of $M_F=yS$ below is the coupling of \eqref{eq:Saction} divided by $\sqrt2$, while the
coupling $g$ is the same as in \cref{thm:compat}. The vector mass matrix is $(M_V^2)_{ab}=2g^2\Re\Tr(T_aS+ST_a^{\!\top})^\dagger(T_bS+ST_b^{\!\top})$
with one coupling for $\su4$ and $\su2_R$, and the eight right-handed Weyl fermions get the Majorana matrix $M_F=yS$.

\begin{theorem}[One-loop selection on $(\mathbf{10},\mathbf1,\mathbf3)$]
\label{thm:oneloop1013}
At leading logarithm, $\Tr M_V^4=g^4\sum_kc_kI_k$ with $c=(33,3,-3,39,5)$ and $\Tr(M_F^\dagger M_F)^2=y^4\sum_ka_kI_k$, and the fluctuations generate the gaps
\[
\Delta_k=\frac L{64\pi^2}\,g^4\bigl(G_k-x\,F_k\bigr),\qquad G=(90,108,-18,84),\quad F=(3,3,\tfrac32,2),\qquad x=\frac{\sum_{\rm fam}y^4}{g^4}.
\]
Gauge fluctuations favour the Standard-Model orbit in the directions $k=2,3,5$, and matter fluctuations disfavour it in every direction. The Standard-Model
orbit is a non-degenerate local minimum on $\mathcal T$ iff $x<30$, the binding conditions being those of $\Delta_{ee}$ and $\Delta_{dd}$. The configuration $\id_4\otimes\id_2$
has excess $\tfrac L{64\pi^2}g^4(48-\tfrac74x)$ over the Standard-Model orbit, negative for $x>\tfrac{192}7\approx27.43$.
\end{theorem}

At $x=0$ the radiatively generated masses are, in units of $\tfrac L{64\pi^2}g^4v^2$, $m_{ee}^2=120$, $m_{ue}^2=84$, $m_{de}^2=144$, $m_{uu}^2=144$, $m_{ud}^2=156$ and
$m_{dd}^2=120$, which satisfy \eqref{eq:sumrules} (v1-12, v1-13); the fifty massive scalars are pseudo-Goldstone bosons of the accidental $O(60)$ symmetry of the tree potential
\cite{Weinberg72,GildenerWeinberg76}. Minimisation on the unit sphere of $\mathcal T$ from random starts finds no configuration below the Standard-Model
orbit at $x=0,20,27$, and at $x=28,29$ finds the orbit of $\id_4\otimes\id_2$ as the lowest end point (v3-2; sixteen archived starts per value, numerical); its excess
$48-\tfrac74x$ crosses zero at $x=192/7$ (v1-14); between $192/7$ and $30$ the Standard-Model vacuum is metastable on $\mathcal T$. We conjecture that on $\mathcal T$ the Standard-Model orbit is the global minimum of the one-loop potential exactly
for $x<192/7$.

With $M_N=yv$ and $M_{W_R}=gv$, $x=\sum_{\rm fam}(M_N/M_{W_R})^4$; in the couplings of \eqref{eq:Saction} and \cref{thm:compat},
$x=\sum_{\rm fam}y^4/(4g^4)$. Stated in terms of masses, $x$ does not depend on the normalisation. The content of the bound is structural: the Majorana Yukawa coupling, matter's only handle on the
vacuum, pushes towards configurations in which the Majorana mass is shared among several fermions, and the gauge fluctuations hold the Standard-Model orbit only
while the heavy neutrinos are not much heavier than $W_R$.

% --- source: tex/s20_selection#6 (The full kernel space)
\subsection{The full kernel space}
\label{sec:fullK}

The general invariant potential on all of $\mathcal M$ has too many couplings to be classified usefully: the numbers of linearly independent Pati--Salam-invariant
real polynomials of degrees $1,2,3,4$ are $0,18,0,686$, and with fermion number imposed $0,12,0,402$. Two structures do not depend on these couplings.

\begin{proposition}[A left--right involution]
\label{prop:LR}
Let $O=E_1E_7E_8$ restricted to $V_*$. It commutes with $P_*$ and with the rotations, satisfies $O^{\!\top}K^0O=K^0$, $O^{\!\top}K^iO=-K^i$, $O\Gamma O^{\!\top}=-\Gamma$ and
$OJO^{\!\top}=J$, and maps $t^L_i$ to $\pm t^R_i$ and each colour bivector to $\pm$ itself. Hence $\psi(t,\mathbf x)\mapsto O\psi(t,-\mathbf x)$, $\Phi\mapsto O\Phi O^{\!\top}$,
with the corresponding map of the gauge fields, is a symmetry of the kinetic, Yukawa and gauge terms of \eqref{eq:Saction}, and of the trace invariants. The image
$\Phi_0'=O\Phi_0O^{\!\top}$ is a kernel in $(\mathbf{10},\mathbf3,\mathbf1)$ orthogonal to $\Phi_0$, the mirror vacuum of left--right models, and along the great circle
$\hat u(\theta)=\cos\theta\,\hat u_0+\sin\theta\,\hat u_0'$, $Q=\tfrac14-\tfrac18\sin^22\theta$ (v2-10).
\end{proposition}

The Standard-Model vacuum and its mirror are exchanged by a symmetry of the action, so an $O$-invariant potential cannot prefer one of them; the choice between
breaking $\su2_R$ and breaking $\su2_L$ is spontaneous, as in models of spontaneous parity violation \cite{SenjanovicMohapatra75}.

Now take the tree potential degenerate on all of $\mathcal M$, $V_0=-\mu^2N+\alpha N^2$, invariant under the accidental $O(272)$. For a unit vector $\hat u\in\mathcal M$ put
\[
\begin{gathered}
Q(\hat u)=\Tr(\hat u^{\!\top}\hat u)^2,\qquad T(\hat u)=\Tr\bigl(G^{-1}K(\hat u)\bigr)^2,\qquad K(\hat u)_{ab}=\langle[t_a,\hat u],[t_b,\hat u]\rangle,\\
E_x(\hat u)=8x\,Q(\hat u)-\tfrac3{16}T(\hat u),
\end{gathered}
\]
with the $21$ Pati--Salam generators $t_a$ and the module metric $G$. The orientation-dependent part of the leading logarithm is then
$\delta V=\tfrac L{64\pi^2}g^4v^4E_x(\Phi/|\Phi|)$ with $v=|\Phi|/\sqrt2$, and at the Standard-Model point $E_x(\hat u_0)=2x-99$.

\begin{theorem}[One-loop second variation on all of $\mathcal M$]
\label{thm:fullsecond}
Let $H(x)$ be the second variation of $E_x$ at $\hat u_0$ on the $271$-dimensional tangent space of the unit sphere.
\begin{enumerate}[label=(\roman*),leftmargin=2.2em]
\item $H(x)=H_G+xH_F$ with $H_F=-8(\id-\Pi_\nu)$, where $\Pi_\nu$ is the orthogonal projector of rank $32$ onto the kernels that pair $\nu_R$ with a Weyl fermion of the
block. $H_G$ preserves each Pati--Salam component of $\mathcal M$.
\item With $m^2=\tfrac12\times$(eigenvalue of $H(x)$) in units of $\tfrac L{64\pi^2}g^4v^2$, the $271$ angular modes have the masses of \cref{tab:271}.
\item $H(x)$ is non-negative with kernel exactly the nine Goldstone directions if and only if $x<27$. At $x=27$ the six modes of $\ell_L\ell_L$, the left-handed triplet,
become massless, and for $x>27$ they are tachyonic.
\end{enumerate}
\end{theorem}

The second variation is computed from the definitions on the $272$-dimensional kernel space, and every entry of \cref{tab:271} is reproduced at
several values of $x$ (v2-4 to v2-8); its restriction to $(\mathbf{10},\mathbf1,\mathbf3)$ reproduces the tensor-model spectrum of \cref{thm:oneloop1013}, which ties
the two normalisations together (v2-9). This is floating-point linear algebra on exactly specified matrices; (iii) follows from the table, whose
lowest line is $108-4x$.

\begin{table}[t]
\centering
\footnotesize
\setlength{\tabcolsep}{4pt}
\begin{tabular}{@{}llccl@{}}
\toprule
fermion pair & multiplet & $\dim_\R$ & component & $m^2$\\
\midrule
$\nu_R\nu_R$ (phase), $\nu_Re_R$, $\nu_Ru_R$ & $(1,1)_0,(1,1)_1,(3,1)_{2/3}$ & $1+2+6$ & $(\mathbf{10},\mathbf1,\mathbf3)$ & $0$ (Goldstone)\\
$e_Re_R$ & $(1,1)_2$ & $2$ & $(\mathbf{10},\mathbf1,\mathbf3)$ & $120-4x$\\
$d_Re_R$ & $(3,1)_{4/3}$ & $6$ & $(\mathbf{10},\mathbf1,\mathbf3)$ & $144-4x$\\
$u_Ru_R$ & $(6,1)_{4/3}$ & $12$ & $(\mathbf{10},\mathbf1,\mathbf3)$ & $144-4x$\\
$u_Rd_R$ & $(6,1)_{1/3}$ & $12$ & $(\mathbf{10},\mathbf1,\mathbf3)$ & $156-4x$\\
$d_Rd_R$ & $(6,1)_{2/3}$ & $12$ & $(\mathbf{10},\mathbf1,\mathbf3)$ & $120-4x$\\
$u_Re_R\pm d_R\nu_R$ & $(3,1)_{1/3}$ & $6+6$ & $(\mathbf{10},\mathbf1,\mathbf3)$, $(\mathbf6,\mathbf1,\mathbf1)_R$ & eigenvalues of $M_3$\\
$u_Rd_R$ (antisymmetric) & $(3,1)_{1/3}$ & $6$ & $(\mathbf6,\mathbf1,\mathbf1)_R$ & $180-4x$\\
$u_Le_L-d_L\nu_L$, $u_Ld_L$ & $(3,1)_{1/3}$ & $6+6$ & $(\mathbf6,\mathbf1,\mathbf1)_L$ & $180-4x$\\
$\ell_L\ell_L$ & $(1,3)_1$ & $6$ & $(\mathbf{10},\mathbf3,\mathbf1)$ & $108-4x$\\
$q_L\ell_L$ & $(3,3)_{1/3}$ & $18$ & $(\mathbf{10},\mathbf3,\mathbf1)$ & $156-4x$\\
$q_Lq_L$ & $(6,3)_{1/3}$ & $36$ & $(\mathbf{10},\mathbf3,\mathbf1)$ & $180-4x$\\
$q_L$--$u_R$, $q_L$--$d_R$ (octet) & $(8,2)_{1/2}$ & $64$ & $(\mathbf{15},\mathbf2,\mathbf2)$ & $174-4x$\\
$\ell_L$--$u_R$, $q_L$--$e_R$ & $(3,2)_{7/6}$ & $24$ & $(\mathbf{15},\mathbf2,\mathbf2)$ & $162-4x$\\
$\ell_L$--$d_R$ & $(3,2)_{1/6}$ & $12$ & $(\mathbf{15},\mathbf2,\mathbf2)$ & $114-4x$\\
$q_L$--$\nu_R$ & $(3,2)_{1/6}$ & $12$ & $(\mathbf{15},\mathbf2,\mathbf2)$ & $114$\\
$q_L$--$d_R$, $\ell_L$--$e_R$ (singlets) & $(1,2)_{1/2}$ & $4+4$ & $(\mathbf{15},\mathbf2,\mathbf2)$, $(\mathbf1,\mathbf2,\mathbf2)$ & $138-4x$, $186-4x$\\
$q_L$--$u_R$, $\ell_L$--$\nu_R$ (singlets) & $(1,2)_{1/2}$ & $4+4$ & $(\mathbf{15},\mathbf2,\mathbf2)$, $(\mathbf1,\mathbf2,\mathbf2)$ & eigenvalues of $M_2$\\
\bottomrule
\end{tabular}
\caption{The radiatively generated masses of the $271$ angular modes of the kernel field at the Standard-Model vacuum, in units of $\tfrac L{64\pi^2}g^4v^2$, with
$x=\sum_{\rm fam}(M_N/M_{W_R})^4$ (\cref{thm:fullsecond}). Pairs written $ff'$ are of Majorana type, pairs written $f$--$f'$ of Dirac type; multiplets are given
under $SU(3)\times SU(2)_L$ with $|Y|$. The two mixing matrices are
$M_3=\bigl(\begin{smallmatrix}84-2x&2x\\2x&180-2x\end{smallmatrix}\bigr)$ and $M_2=\bigl(\begin{smallmatrix}138-x&\sqrt3x\\\sqrt3x&186-3x\end{smallmatrix}\bigr)$, with eigenvalues
$132-2x\mp2\sqrt{x^2+576}$ and $162-2x\mp2\sqrt{x^2-12x+144}$.}
\label{tab:271}
\end{table}

Statement (i) is the mechanism by which matter opposes the Standard-Model orbit, now on all of $\mathcal M$. At fixed $|\Phi|$ the fermion loop $2\sum m_f^4$ is largest when the
whole Majorana mass sits on $\nu_R$; it decreases quadratically along every direction that moves mass onto pairs not containing $\nu_R$ and is stationary along the
directions that pair $\nu_R$ with another fermion. Every mode that does not involve $\nu_R$ therefore has slope $-4$, and the first of them to become massless is the one
with the smallest gauge contribution, the left-handed triplet. The mixing in $M_3$, between the $(\mathbf{10},\mathbf1,\mathbf3)$ mode $u_Re_R+d_R\nu_R$ and the $(\mathbf6,\mathbf1,\mathbf1)_R$
mode $u_Re_R-d_R\nu_R$, is an instance of the mixing announced at the end of \cref{sec:kernels}: it moves the $\Delta_{ue}$ threshold from $42$ to $315/11$. The
thresholds are the zeros of the lowest lines,
\[
x=27\ (\ell_L\ell_L),\qquad x=28.5\ (\ell_L\text{--}d_R),\qquad x=\tfrac{315}{11}\approx28.64\ (M_3),\qquad x=30\ (e_Re_R,\ d_Rd_R),
\]
and \cref{fig:radiative} shows them. Along the great circle of \cref{prop:LR},
\[
E_x(\hat u(\theta))=2x-99+(27-x)\sin^22\theta ,
\]
so for $x>27$ the Standard-Model vacuum descends towards the left--right symmetric configuration $\theta=\pi/4$, where $\nu_R$ and a left-handed lepton carry equal
Majorana masses.

\begin{figure}[t]
\centering
\begin{tikzpicture}
\begin{axis}[width=12.5cm,height=7.6cm,xmin=0,xmax=32,ymin=-15,ymax=200,
  xlabel={$x=\sum_{\rm fam}(M_N/M_{W_R})^4$},ylabel={$m^2$ in units of $\frac{L}{64\pi^2}g^4v^2$},
  axis lines=left,clip=true,xtick={0,5,10,15,20,25,30},ytick={0,50,100,150,200},
  ylabel style={font=\small},xlabel style={font=\small},
  legend style={at={(0.02,0.03)},anchor=south west,font=\footnotesize,fill=white,draw=mgray},legend cell align=left]
\fill[lgray!60] (axis cs:24.884,-15) rectangle (axis cs:27,200);
\addplot[gray!70,thin,forget plot,domain=0:32] {144-4*x};
\addplot[gray!70,thin,forget plot,domain=0:32] {156-4*x};
\addplot[gray!70,thin,forget plot,domain=0:32] {162-4*x};
\addplot[gray!70,thin,forget plot,domain=0:32] {174-4*x};
\addplot[gray!70,thin,forget plot,domain=0:32] {180-4*x};
\addplot[gray!70,thin,forget plot,domain=0:32] {186-4*x};
\addplot[gray!70,thin,forget plot,domain=0:32] {114};
\addplot[gray!70,thin,forget plot,domain=0:32,samples=60] {162-2*x+2*sqrt(x^2-12*x+144)};
\addplot[gray!70,thin,forget plot,domain=0:32,samples=60] {162-2*x-2*sqrt(x^2-12*x+144)};
\addplot[gray!70,thin,forget plot,domain=0:32,samples=60] {132-2*x+2*sqrt(x^2+576)};
\addplot[scol,very thick,domain=0:32] {108-4*x}; \addlegendentry{$\ell_L\ell_L$: $108-4x$}
\addplot[ocol,thick,domain=0:32] {114-4*x}; \addlegendentry{$\ell_L$--$d_R$: $114-4x$}
\addplot[gcol,thick,domain=0:32,samples=80] {132-2*x-2*sqrt(x^2+576)}; \addlegendentry{lower eigenvalue of $M_3$}
\addplot[tcol,thick,domain=0:32] {120-4*x}; \addlegendentry{$e_Re_R$, $d_Rd_R$: $120-4x$}
\draw[black,dashed] (axis cs:0,0) -- (axis cs:32,0);
\draw[gray,densely dotted] (axis cs:27,-15) -- (axis cs:27,200);
\draw[gray,densely dotted] (axis cs:28.5,-15) -- (axis cs:28.5,200);
\draw[gray,densely dotted] (axis cs:28.636,-15) -- (axis cs:28.636,200);
\draw[gray,densely dotted] (axis cs:30,-15) -- (axis cs:30,200);
\node[anchor=north east,font=\scriptsize,align=right] at (axis cs:24.7,196) {metastable band\\ $x_*\le x\le27$};
\end{axis}
\end{tikzpicture}
\caption{Radiatively generated masses of the angular modes of the kernel field at the Standard-Model vacuum (\cref{thm:fullsecond}, \cref{tab:271}). The coloured lines
are the four lowest; they cross zero at $x=27$, $28.5$, $315/11$ and $30$ (dotted). The Standard-Model orbit is a non-degenerate local minimum for $x<27$. That it is the
global minimum for $x<x_*\approx24.884$, and metastable in the shaded band, is conjectured from the searches of \cref{prop:randomstart}.}
\label{fig:radiative}
\end{figure}

\begin{proposition}[Random-start minimisation]
\label{prop:randomstart}
$E_x$ was minimised on the unit sphere of $\mathcal M$ by a quasi-Newton method (L-BFGS) from standard Gaussian starts. For $x=0,20,22$ no run (forty per value)
ended below the energy $2x-99$ of the Standard-Model orbit, and at least $39$ of the $40$ ended at that orbit or its mirror; which end points above it
are reached depends on the floating-point path (in the regressions of the release candidate and of version 1.0.0, with one BLAS thread per
script, one run at $x=22$ stopped at the higher critical value $E=-51.562388$). At $x=25$, about one run in five ($22$ or $23$ of $100$, according to the floating-point environment) ended below it, at $E=-49.1398704$, in two forms exchanged by
the involution of \cref{prop:LR}. The form with $\su2_L$ unbroken has weight $0.968$ on $(\mathbf{10},\mathbf1,\mathbf3)$ and $0.032$ on $2(\mathbf6,\mathbf1,\mathbf1)$, and its stabiliser
among the $21$ Pati--Salam generators has dimension nine, contains $\su2_L$ and projects onto colour with rank six, whereas any algebra containing $\su3$ projects with
rank at least eight: colour is broken. There four right-handed Weyl fermions carry Majorana masses, one of them $2.04$ times as heavy as the other three. A second
optimiser (nonlinear conjugate gradient) reaches the same energy. Followed in $x$, this minimum crosses the energy of the Standard-Model orbit at
$x_*=24.88405$ (to the stated digits), where it is a local minimum on the sphere.
\end{proposition}

The proposition reports numerical experiments; it is not a certificate. They are archived with their seeds, start distribution, optimisers, tolerances and
residuals in the third script (v3) of \path{verification/ch07_mass_vacuum}, checks v3-3 to v3-7: the largest final Riemannian gradient is below $10^{-5}$, the Goldstone zero modes are
left to the optimiser (the energies of the end points agree to $10^{-9}$ between runs and optimisers), and $x_*$ is located by Brent's method on warm-started
minimisations with tolerance $10^{-10}$, giving $24.8840501$. The fraction of runs that reaches the competitor depends on the start distribution; with standard Gaussian starts it is $22\%$. At $x=24.5$ no run of forty ends below the Standard-Model energy. Only the local statement,
\cref{thm:fullsecond}, is a theorem about the stated approximation; a random search does not establish a global minimum.

We conjecture that in the radiative scenario on $\mathcal M$ the Standard-Model orbit and its mirror are the global minima exactly for $x<x_*$, and that for $x_*<x<27$ the
Standard-Model vacuum is metastable. The competitor again spreads the Majorana mass, now over four right-handed fermions instead of one, which is what the fermion loop
rewards, at the price of breaking colour, which the gauge loop penalises. In terms of masses the local condition and the conjectured global one read
\[
\begin{aligned}
&M_N<27^{1/4}M_{W_R}=2.280\,M_{W_R}\quad\text{and}\quad M_N<2.233\,M_{W_R}&&\text{(one family dominating)},\\
&M_N<9^{1/4}M_{W_R}=1.732\,M_{W_R}\quad\text{and}\quad M_N<1.697\,M_{W_R}&&\text{(three degenerate families)}.
\end{aligned}
\]
They hold under the following restrictions: a tree potential degenerate on $\mathcal M$ at the matching scale, one loop at leading logarithm, $\kappa=0$, which fermion number
makes automatic, a common gauge coupling for $\su4$ and $\su2_R$, and a single Yukawa coupling for all components of the kernel field (\cref{ass:field}). The masses of
\cref{tab:271} are curvatures of the leading-logarithm effective potential divided by the kinetic norm; pole masses would need the wave-function renormalisation and a
gauge-consistent perturbative expansion, which are not attempted. The bounds are therefore statements about this model and approximation, not bounds on heavy-neutrino
masses outside it. Within it, the result is that gauge fluctuations make the Standard-Model orientation of the
Majorana condensate stable in every direction of the kernel space as long as the heavy neutrinos are not much heavier than $W_R$; beyond that, the local instability is
the condensation of the left-handed triplet, and the configuration that overtakes the Standard-Model orbit globally, at the slightly smaller $x_*$, breaks colour.

% --- source: tex/s21_matter#0 ((head))


% --- source: tex/s21_matter#1 (Matter in the selected vacuum)
\section{Matter in the selected vacuum}
\label{sec:matter}

We return to the vacuum $\Phi=r\Phi_0$, $r\neq0$, with $v=|\Phi|/\sqrt2=\sqrt2\,r$, and ask whether the matter of the block is compatible with
it in the sense stated at the beginning of \cref{sec:kernels}.

\begin{theorem}[Compatibility with matter]
\label{thm:compat}
\begin{enumerate}[label=(\alph*),leftmargin=2.2em]
\item \emph{Unbroken algebra.} Among the $21$ Pati--Salam generators the stabiliser of $\Phi_0$ has dimension $12$ and closes under brackets; its part
orthogonal to $\su3\oplus\su2_L$ is spanned by $t_Y=t^R_3+\tfrac16H$. Among the $23$ generators it has dimension $14$; the two further directions $J+\tfrac13H$
and $J-J\Gamma$ carry the charges $4B$ and $1-\Gamma$ and are global.
\item \emph{Vector bosons.} The matrix $M_{ab}=\langle[t_a,\Phi_0],[t_b,\Phi_0]\rangle/|\Phi_0|^2$, taken relative to the module metric, has eigenvalues $0$ (twelve
times), $4$ (eight times) and $20$ (once), so the squared vector masses are $0$, $g^2v^2$ and $5g^2v^2$ in the conventionally normalised coupling $g=\sqrt8\,g_A$ of
\cref{sec:kernels}, that is $8g_A^2v^2$ and $40g_A^2v^2$ for the module-normalised fields of \eqref{eq:Saction}. The eight carry $Y=\pm\tfrac23$ (six leptoquarks
$X$) and $\pm1$ (the two $W_R^\pm$); the one at $5g^2v^2$ is neutral, a $Z'$.
\item \emph{Couplings.} In the module trace $|t^L_3|^2=\tfrac18$ and $|t_Y|^2=\tfrac5{24}$, so
\[
\sin^2\theta_W=\frac{|t^L_3|^2}{|t^L_3|^2+|t_Y|^2}=\frac38
\]
at the matching scale, the value of the Pati--Salam and $\so{10}$ embeddings \cite{GeorgiQuinnWeinberg74}.
\item \emph{Dirac operator.} The kinetic kernels are symmetric with $K^0=\id$ and are preserved by the twelve unbroken generators, and $H(p)^2=|p|^2+y^2r^2P_\nu$.
\item \emph{Chirality.} The massless fermions are $(3,2)_{1/6}\oplus(1,2)_{-1/2}$ with $\Gamma=-1$ and $(3,1)_{2/3}\oplus(3,1)_{-1/3}\oplus(1,1)_{-1}$ with $\Gamma=+1$:
the unbroken $\su2_L$ acts on one Lorentz chirality only.
\item \emph{Anomalies.} On the block the $Y^3$, mixed gravitational, $SU(2)^2Y$, $SU(3)^2Y$ and $SU(3)^3$ anomalies vanish, with and without $\nu_R$, and there are
four left-handed $SU(2)$ doublets, an even number \cite{Witten82}.
\end{enumerate}
\end{theorem}

The formula in (c) deserves a line. With one coupling for the whole Pati--Salam algebra, the effective couplings of $t^L_3$ and $t_Y$ are in the ratio of
their norms in the module, $g'^2/g^2=|t^L_3|^2/|t_Y|^2=\tfrac{3/24}{5/24}=\tfrac35$, the familiar normalisation of hypercharge in unified models, and
$\sin^2\theta_W=g'^2/(g^2+g'^2)=\tfrac38$, at the (unspecified) scale where the single coupling holds; without a matching scale, spectrum, running and
thresholds this is a normalisation statement, not a prediction of the low-energy angle. Every requirement of a matter-compatible vacuum is met at this stage except mass
generation for the charged fermions, which needs the electroweak stage.

% --- source: tex/s21_matter#2 (The electroweak stage and the fermion masses)
\subsection{The electroweak stage and the fermion masses}

Everything in this subsection assumes, in addition, a vacuum value of the bidoublet kernels; their dynamics and their scale are not derived. In the radiative
scenario of \cref{sec:fullK} the colour-singlet doublets have squared masses $138-4x$, $186-4x$ and the eigenvalues of $M_2$, so they sit at the Pati--Salam
scale and a light doublet requires a tuning that we do not derive.

\begin{proposition}[Electroweak stage]
\label{prop:EW}
\begin{enumerate}[label=(\roman*),leftmargin=2.2em]
\item The kernels commuting with $\su3$ and with $Q=T^L_3+Y$ form a $12$-dimensional space: two each in $(\mathbf{10},\mathbf1,\mathbf3)$ and $(\mathbf{10},\mathbf3,\mathbf1)$ and four
each in $(\mathbf{15},\mathbf2,\mathbf2)$ and $(\mathbf1,\mathbf2,\mathbf2)$. In each bidoublet the up-type kernels (no $d,e$ entries) and the down-type kernels (no $u,\nu$ entries)
form two-dimensional spaces.
\item On the charge sectors $u$ ($Q=\tfrac23$), $d$ ($-\tfrac13$), $\nu$ ($0$) and $e$ ($-1$), the $(\mathbf1,\mathbf2,\mathbf2)$ kernels give $m_u=m_\nu^D$ and $m_d=m_e$; the
$(\mathbf{15},\mathbf2,\mathbf2)$ kernels give lepton masses three times the quark masses, with relative sign $-3$ once aligned with a $(\mathbf1,\mathbf2,\mathbf2)$ kernel. For
$\Phi=\kappa_1K_1+\kappa_{15}K_{15}$ of each type,
\[
m_{\rm quark}=|\kappa_1+\kappa_{15}|,\qquad m_{\rm lepton}=|\kappa_1-3\kappa_{15}|,
\]
the Georgi--Jarlskog factor \cite{GeorgiJarlskog79} used in fits of fermion masses with $\mathbf{10}\oplus\mathbf{126}$ \cite{BabuMohapatra93}.
\item With $\Phi=v_R\Phi_0+\kappa K_{1u}$ the neutral sector has singular values $\tfrac12\bigl(\sqrt{v_R^2+4\kappa^2}\mp v_R\bigr)$, each four times: the seesaw pair
\cite{Minkowski77,MohapatraSenjanovic80}, with light mass $\approx\kappa^2/v_R$.
\item For $\Phi=v_R\Phi_0+v_uK_{1u}+v_dK_{1d}$ the stabiliser among the $21$ generators is nine-dimensional, $\su3\oplus\uu1_Q$, and $Q$ annihilates $\Phi$.
\end{enumerate}
\end{proposition}

\noindent Here $K_{1u}$ and $K_{1d}$ are up- and down-type $(\mathbf1,\mathbf2,\mathbf2)$ kernels of (i), normalised to $|K|^2=n_K=32$; each is fixed up to a
rotation in its two-dimensional space, which changes none of the statements. Executable check: a8-2, a8-3, a8-4, a8-5 (numerical; \cref{app:methods}).

\begin{example}[An electroweak spectrum]
\label{ex:EWspectrum}
At $(v_R,v_u,v_d)=(20,1,\tfrac12)$ the squared vector masses, computed from \eqref{eq:Saction} as written, that is for $g_A=1$ and $g^2=8$ (divide
by $8$ for $g=1$), are $0$ (nine times), $79.360$ (twice, $W^\pm$),
$127.872$ (once, $Z$), $6400$ (six times), $6480.64$ (twice) and $32032.1$ (once). The ratio $m_W^2/m_Z^2=0.62062$ is close to $\cos^2\theta_W=\tfrac58$, and at
$v_R=200$ it differs from $\tfrac58$ by $-4.4\times10^{-5}$: $\rho=1$ at tree level up to corrections of order $\kappa^2/v_R^2$. In the same vacuum at
$(v_R,v_u,v_d)=(20,3,1)$ the fermion mass eigenspaces are a light Majorana neutrino of mass $(\sqrt{436}-20)/2\approx0.440$ and a heavy one of mass
$(\sqrt{436}+20)/2\approx20.440$, each of real dimension four, the $d$ quarks and the electron with $m_d=m_e=1$, and the $u$ quarks with mass $3$. The
chirality content of the neutral pair is $\langle\Gamma\rangle=\mp\cos2\vartheta$ with $\tan2\vartheta=2v_u/v_R$, which is their mixing, and the Majorana-type part of the full
mass operator is exactly $v_R\Phi_0$: $\nu_L$ acquires its Majorana mass only through the mixing.
\end{example}

What is not obtained here is equally definite: the values of $v_u$ and $v_d$, the potential that would select them together with $v_R$, and the
hierarchy $\kappa\ll v_R$, which in Pati--Salam and left--right models is a tuning of the scalar potential \cite{Gildener76}.

% --- source: tex/s21_matter#3 (What the vacuum does not supply)
\subsection{What the vacuum does not supply}

The statements of this subsection are impossibility results within explicitly stated algebras of operators. They say where an input enters, not that no
mechanism outside those algebras could replace it.

\begin{proposition}[Chirality is field content]
\label{prop:mirror}
Let $\mathcal A$ be the algebra of operators on $S_3$ generated by the $91$ internal bivectors and $J$. Then $\gamma^0$ commutes with every element of $\mathcal A$ and
with the rotations, anticommutes with $\Gamma$ and the boosts, and maps $V_*$ isometrically onto the orthogonal rank-$64$ subspace $\gamma^0V_*$, on which $\mathcal A$
acts by the same matrices and $\Gamma$ has the opposite sign. Hence every eigenspace of an element of $\mathcal A$ has a mirror image with the same eigenvalue and the
opposite chirality, and no vacuum value of a field with values in $\mathcal A$, nor any function of such values, separates $V_*$ from $\gamma^0V_*$.
\end{proposition}

The block is selected by $\mathcal X=\Gamma\chi_W$, which contains the spacetime orientation through $\Gamma=J\gamma^0\gamma^1\gamma^2\gamma^3$ and is not in $\mathcal A$. The kernel field
is defined on $V_*$ and presupposes it, so it does not escape the proposition either. The Standard-Model vacuum of \cref{thm:select0} is chiral because the
matter on which it acts already is; this is the arena's form of the fact that a vector-like gauge theory does not generate chirality through its vacuum
\cite{VafaWitten84b}.

\begin{proposition}[Hermitian order parameters keep rank five]
\label{prop:rankfive}
On the spin-up part of the block, the operators commuting with $J$, $\Gamma$ and the twelve Standard-Model generators form a twelve-dimensional algebra whose
symmetric part is six-dimensional: the commutative algebra of the six multiplet projectors $(Q,L,u,d,e,\nu)$, which are pairwise inequivalent. A Hermitian vacuum
value in it commutes with $T^R_3$ and $\BmL$, so its stabiliser contains $\su3\oplus\su2_L\oplus\uu1_{T^R_3}\oplus\uu1_{\BmL}$, of rank five; the Standard-Model algebra, of
rank four, is not the unbroken algebra of such an order parameter.
\end{proposition}

\noindent The spin-up part is an eigenspace of $Je_4e_5$ on the block. Executable check: a8-7 (numerical; \cref{app:methods}), which also finds the six
multiplet projectors with multiplicities $12,4,6,6,2,2$.

This is the block version of the statement that order parameters of adjoint type do not lower the rank \cite{Li74}. The kernel field escapes it because a
Grassmann bilinear carries charge: $[t^R_3,\Phi_0]=J\Phi_0\neq0$.

\begin{proposition}[Gauge alignment does not see chirality]
\label{prop:alignment}
Consider rank-$64$ projectors $P$ commuting with Lorentz and $J$, with the alignment energy $V(P)=\tfrac1{512}\sum_T\|[T,P]\|^2$ over the $43$ compact internal
bivectors. For $P_*$, $\Pfam$, the colour-singlet lepton projector $P_{\rm lep}$ and the vector-like $Q_v=P_L+\gamma^0P_L\gamma^0$ with $P_L=P_*\tfrac12(\id-\Gamma)$:
\begin{center}
\small
\begin{tabular}{@{}lcccc@{}}
\toprule
 & $P_*$ & $\Pfam$ & $P_{\rm lep}$ & $Q_v$\\
\midrule
$V(P)$ & $5$ & $9$ & $3$ & $5$\\
stabiliser among the $43$ & $23$ & $9$ & $37$ & $25$\\
axial, $P\gamma^0P=0$ & yes & --- & --- & no\\
\bottomrule
\end{tabular}
\end{center}
Alignment prefers $P_{\rm lep}$ to $P_*$, and $V(P_*)=V(P_L)+V(P_R)$ with $V(P_L)=V(P_R)=\tfrac52$ while $V(Q_v)=2V(P_L)=V(P_*)$: the alignment energy does not see
chirality.
\end{proposition}

Vacuum alignment \cite{Peskin80,Preskill81} selects among projectors only by the amount of gauge symmetry they keep; it neither selects the Pati--Salam block nor
distinguishes a chiral block from a vector-like one of the same rank. Finally, with the single Yukawa coupling of \cref{ass:field}, matter opposes the
Standard-Model vacuum: on $(\mathbf{10},\mathbf1,\mathbf3)$ every $F_k$
of \cref{thm:oneloop1013} is positive, and on all of $\mathcal M$ the Yukawa part of the second variation is $-8x(\id-\Pi_\nu)\le0$. The fermion loop lowers the energy
along every direction that moves Majorana mass away from $\nu_R$ and raises it along none; on its own it would select configurations in which the mass is
spread, which break colour. The matter sector is not the agent that selects the Standard-Model orbit. At one loop the gauge fluctuations are, and they succeed
for $x<27$ locally (\cref{thm:fullsecond}), and globally, by conjecture from numerical searches, for $x<x_*$ (\cref{prop:randomstart}).

% --- source: tex/s22_goldstone#0 ((head))


% --- source: tex/s22_goldstone#1 (Goldstone bosons, antiparticles, baryon number and topology)
\section{Goldstone bosons, antiparticles, baryon number and topology}
\label{sec:goldstone}

The selected vacuum can be examined from four further angles: which of its zero modes are Goldstone bosons and which are eaten; how the discrete
operations that exchange particles and antiparticles act on it; whether the potential protects baryon number; and what topological defects its orbit
supports. Each answer depends on counting the two global central generators $J$ and $J\Gamma$ alongside the $21$ gauged ones.

% --- source: tex/s22_goldstone#2 (Goldstone bosons)
\subsection{Goldstone bosons}

\begin{proposition}[Goldstone directions attached to the symmetries]
\label{prop:global}
$J$ and $J\Gamma$ are antisymmetric, commute with the kinetic kernels and the $21$ Pati--Salam generators, span the centre of the $23$-dimensional faithful
algebra, are anomalous, and carry no gauge fields. If $V$ is invariant under all $23$ generators and $\Phi$ is stationary, the tangent space of the $23$-orbit
lies in the kernel of the Hessian; the gauge bosons eat the tangents of the $21$-orbit, whose dimension is the rank of the vector mass matrix, and
\[
\#\{\text{symmetry-generated Goldstone directions not eaten}\}=\dim(23\text{-orbit})-\dim(21\text{-orbit}).
\]
\end{proposition}

The count is classical, and it counts the Goldstone directions attached to the $23$ specified generators, not all massless scalars: a potential with an
accidental symmetry has more zero modes, as the degenerate tree potential of \cref{sec:fullK} has $271$ against a gauge orbit of dimension nine. Since $J$
and $J\Gamma$ are anomalous, an uneaten direction is a pseudo-Goldstone boson, massive once the anomaly is accounted for.

\begin{theorem}[No Majoron]
\label{thm:nomajoron}
At $r\Phi_0$ the $21$-orbit and the $23$-orbit both have dimension nine, so the global symmetries add no Goldstone direction beyond those eaten. The reason is
\[
[J,\Phi_0]=[J\Gamma,\Phi_0]=2\,[t^R_3,\Phi_0]:
\]
the fermion-number rotation of the $\nu_R\nu_R$ condensate is a gauge rotation, and $J+\tfrac13H$, of charge $4B$, stays unbroken. The nine eaten modes are the
$\nu_R\nu_R$ phase ($Z'$), $\nu_Re_R$ ($W_R^\pm$) and $u_R\nu_R$ (six leptoquarks), matching the vector masses of \cref{thm:compat}.
\end{theorem}

The would-be Majoron of a spontaneously broken lepton number is eaten by the $Z'$. At the electroweak stage the situation changes.

\begin{theorem}[One uneaten direction at the electroweak vacuum]
\label{thm:axion}
For $\Phi=v_R\Phi_0+v_uK_{1u}+v_dK_{1d}$:
\begin{center}
\small
\begin{tabular}{@{}lccccc@{}}
\toprule
vacuum & stabiliser in $21$ & stabiliser in $23$ & massive vectors & $23$-orbit & Goldstone directions\\
\midrule
generic, $v_uv_d\neq0$ & $9$ & $10$ & $12$ & $13$ & $\mathbf1$\\
$v_d=0$ or $v_u=0$ & $9$ & $11$ & $12$ & $12$ & $0$\\
$v_R=0$ & $16$ & $17$ & $5$ & $6$ & $1$\\
\bottomrule
\end{tabular}
\end{center}
At the electroweak vacuum the tangents of $J$ and $J+\tfrac13H$ lie in the gauge orbit and those of $J\Gamma$ and $J-J\Gamma$ do not; $J+\tfrac13H$ annihilates $\Phi$
and $J-J\Gamma$, of charge $1-\Gamma$, does not.
\end{theorem}

The mechanism is simple. $J\Gamma$ acts on every Dirac-type kernel as the complex structure $2J\Gamma$, so it rotates the up-type and the down-type lines by the
same phase, while the only Standard-Model generator that moves within these lines, the $Z$ direction, rotates them by opposite phases. When both values are
nonzero the $J\Gamma$ direction cannot be absorbed. The broken global charge is $1-\Gamma$, whose anomalies per Weyl fermion in the all-left convention are
\begin{center}
\small
\begin{tabular}{@{}lcccccc@{}}
\toprule
charge & $SU(3)^2$ & $SU(2)_L^2$ & $Y^2$ & $Q^2$ & gravitational & electroweak vacuum\\
\midrule
$1-\Gamma$ & $\mathbf2$ & $4$ & $\tfrac43$ & $\tfrac{16}3$ & $16$ & broken\\
$4B=1+(\BmL)$ & $0$ & $2$ & $-2$ & --- & $0$ & unbroken\\
\bottomrule
\end{tabular}
\end{center}
The current of the broken charge has a colour anomaly, so its classical Goldstone direction is a pseudo-Goldstone boson of Peccei--Quinn--Weinberg--Wilczek
type \cite{PecceiQuinn77,Weinberg78,Wilczek78}: an axion. Since its only order parameter is the bidoublet, its decay constant is of the electroweak scale, and such an
axion is excluded experimentally \cite{PDG}. Within this field content, leaving $U(1)_{1-\Gamma}$ intact is therefore a phenomenological obstruction, and the
electroweak-stage potential must break it explicitly. The trace potential and the one-loop gauge and Yukawa potential of \cref{sec:selection} are invariant
under every orthogonal conjugation commuting with the Pati--Salam algebra, hence under $\exp\theta(J-J\Gamma)$, and the holomorphic quartic lives on the
$\Gamma=+1$ half, where $J-J\Gamma=0$; none of them lifts the axion.

\begin{proposition}[Bidoublet mass terms and their $CP$ parity]
\label{prop:bidoublet}
On $(\mathbf1,\mathbf2,\mathbf2)$ there are three Pati--Salam- and $J$-invariant symmetric forms, one of them $J\Gamma$-invariant. The $CP$ operator of
\cref{thm:fourclasses} acts on them with eigenvalues $(+1,+1,-1)$: the invariant form $B_{\rm inv}$ is $CP$-even, and the two-dimensional breaking space splits into a
$CP$-even $B_e$ and a $CP$-odd $B_o$. Along the $J\Gamma$-orbit,
\[
B_e\bigl(e^{\theta\ad J\Gamma}\Phi\bigr)=-v_uv_d\cos4\theta,\qquad B_o\bigl(e^{\theta\ad J\Gamma}\Phi\bigr)=-v_uv_d\sin4\theta,\qquad B_{\rm inv}=\text{const},
\]
the analogues of $\Re$ and $\Im\Tr\tilde\Phi^\dagger\Phi$ in left--right models.
\end{proposition}

\noindent Executable check: a8-8 (numerical; \cref{app:methods}): the three forms, the $CP$ eigenvalues, and the dependence $A\cos(4\theta+\varphi)$,
$A\sin(4\theta+\varphi)$ along the orbit with $A$ proportional to $v_uv_d$ and $B_{\rm inv}$ constant. The amplitude $-v_uv_d$ and $\varphi=\pi$ of the
display are a choice of normalisation of $B_e$, $B_o$ and of the relative phase of $K_{1u}$ and $K_{1d}$.

\begin{theorem}[The kinetic norm of the axion direction]
\label{thm:mA}
Let $a=J\Phi_0$, $e_u=J\Gamma K_{1u}$, $e_d=J\Gamma K_{1d}$; these are mutually orthogonal with $|a|^2=4$ and $|e_u|^2=|e_d|^2=n_K=32$, and
$t=[J\Gamma,\Phi]=2v_Ra+2v_ue_u+2v_de_d$. The gauge tangent space splits orthogonally into its part in $\operatorname{span}\{a,e_u,e_d\}$, spanned by $a$ and
$z=v_ue_u-v_de_d$, and its part orthogonal to that span, and for all $v_R,v_u,v_d$
\[
\bigl|P_\perp t\bigr|^2=|t|^2-\frac{\langle t,a\rangle^2}{|a|^2}-\frac{\langle t,z\rangle^2}{|z|^2}=\frac{16\,n_K\,v_u^2v_d^2}{v_u^2+v_d^2}.
\]
\end{theorem}

\begin{proof}
$J$ and $J\Gamma$ are orthogonal, so $|a|^2=|\Phi_0|^2=4$ and $|e_{u,d}|^2=|K_{1u,1d}|^2=n_K$. The three vectors are orthogonal because $a$ lies in
$(\mathbf{10},\mathbf1,\mathbf3)$ and $e_u,e_d$ in $(\mathbf1,\mathbf2,\mathbf2)$, and the up- and down-type kernels have disjoint charge sectors. $\Phi_0$
anticommutes with $J$ and is supported on $\Gamma=+1$, and $K_{1u},K_{1d}$ anticommute with $J\Gamma$, which gives the expression for $t$. These facts and
the orthogonal splitting of the gauge tangent space, with the part in the span equal to $\operatorname{span}\{a,z\}$, are recomputed in a8-6 (numerical;
\cref{app:methods}). Given them, $t$ lies in the span, so its projection onto the gauge tangent space is its projection onto $\operatorname{span}\{a,z\}$,
with $a\perp z$; and $\langle t,a\rangle=8v_R$, $\langle t,z\rangle=2n_K(v_u^2-v_d^2)$, $|z|^2=n_K(v_u^2+v_d^2)$, $|t|^2=16v_R^2+4n_K(v_u^2+v_d^2)$, so
\[
\bigl|P_\perp t\bigr|^2=4n_K(v_u^2+v_d^2)-\frac{4n_K(v_u^2-v_d^2)^2}{v_u^2+v_d^2}=\frac{16\,n_K\,v_u^2v_d^2}{v_u^2+v_d^2}.\qedhere
\]
\end{proof}

The heavy-condensate component of $t$ is a gauge direction, which is why $v_R$ drops out. If $\mu_e^2B_e$ is the only term breaking $U(1)_{1-\Gamma}$, with canonical
amplitudes $h_{u,d}=\sqrt{32}\,\kappa_{u,d}$ and $m_{12}^2=\mu_e^2/32$, then to leading order in the explicit breaking
\[
m_A^2=\frac{\partial_\theta^2V}{|P_\perp[J\Gamma,\Phi]|^2}=\frac{m_{12}^2}{\sin\beta\cos\beta},\qquad \tan\beta=\frac{v_u}{v_d},
\]
the pseudoscalar mass of the two-Higgs-doublet model with a soft $m_{12}^2$ term \cite{Branco12}. Phase-dependent quartics add to it, and whether this curvature
is a mass eigenvalue once mixing with the other neutral scalars is included requires the complete Hessian of a specified potential. With $CP$ imposed
($\mu_o^2=0$) every extremum of the breaking mass terms has $\sin4\theta=0$, so no quadratic term can break $CP$ spontaneously; that needs terms of charge $\pm8$
under $J\Gamma$, and the potential has $31+31$ of them among its quartics (a6-5), the analogues of the phase-dependent quartics of left--right models \cite{DGKO91}.

\begin{figure}[t]
\centering
\begin{tikzpicture}[>=Stealth,font=\small,
  stage/.style={draw,thick,rounded corners=4pt,align=center,inner sep=5pt,text width=4.2cm},
  stp/.style={align=left,text width=5.3cm,font=\footnotesize},
  stpl/.style={align=right,font=\footnotesize}]
\node[stage,draw=tcol,fill=tcol!6] (ps) at (0,0) {$\su4\oplus\su2_L\oplus\su2_R$\\ $21$ gauged; global $J$, $J\Gamma$};
\node[stage,draw=gcol,fill=gcol!6] (sm) at (0,-2.6) {$\su3\oplus\su2_L\oplus\uu1_Y$\\ $12$ gauged; global $4B$, $1-\Gamma$};
\node[stage,draw=scol,fill=scol!6] (em) at (0,-5.2) {$\su3\oplus\uu1_Q$\\ $9$ gauged; global $4B$};
\draw[->,thick] (ps) -- (sm); \draw[->,thick] (sm) -- (em);
\node[stp,anchor=west] at (2.6,-1.3) {$v_R\Phi_0$ ($\nu_R\nu_R$ in $(\mathbf{10},\mathbf1,\mathbf3)$)\\ eaten: $Z'$, $W_R^\pm$, $6\,X$ (nine)\\ no Majoron: $[J,\Phi_0]=2[t^R_3,\Phi_0]$\\ $\Z_2$ strings and monopoles, conditionally};
\node[stp,anchor=west] at (2.6,-3.9) {$v_uK_{1u}+v_dK_{1d}$ (bidoublet)\\ eaten: $W^\pm$, $Z$ (three more)\\ one axion-like direction ($1-\Gamma$)\\ unless $V$ breaks $U(1)_{1-\Gamma}$};
\node[stpl,anchor=east] at (-2.6,-1.3) {$\sin^2\theta_W=\tfrac38$\\ $\nu_R$: Majorana, mass $|y|r$};
\node[stpl,anchor=east] at (-2.6,-3.9) {seesaw pair; $m_d=m_e$, $m_u=m_\nu^D$\\ Georgi--Jarlskog factor $-3$};
\end{tikzpicture}
\caption{The two stages of symmetry breaking in the selected vacuum (\cref{sec:matter,sec:goldstone}). Each arrow lists the order parameter, the vectors that
eat the Goldstone directions, and the uneaten directions attached to the global central charges.}
\label{fig:breaking}
\end{figure}

% --- source: tex/s22_goldstone#3 (Particles and antiparticles)
\subsection{Particles and antiparticles}

At one constituent charge conjugation is unique up to a $U(1)_J$ rotation, lies in the surplus plane and is the Majorana real structure (\cref{thm:C}), and the
Fock space of one momentum has a unique charge-zero ground state of energy $-2E$ (\cref{thm:spectra}). That value is the zero-point energy of two particle and two
antiparticle modes; summed over momenta it is the fermion term $-2\sum_{\rm Weyl}$ of the one-loop potential \eqref{eq:V1}, the very term through which matter
opposes the Standard-Model vacuum. At three constituents the discrete operations are classified as follows. A map $\psi\mapsto S\psi$ with $S$ orthogonal on $V_*$
exchanges particles and antiparticles when $SJ=-JS$; it keeps or reverses chirality according to $S\Gamma S^{-1}=\pm\Gamma$; it goes with $\mathbf x\mapsto-\mathbf x$ when
$S^{\!\top}K^iS=-K^i$.

\begin{theorem}[Four classes of discrete operators]
\label{thm:fourclasses}
The blades of $\Cl(9,9)$ that commute with the three involutions of $P_*$ and with the rotations number $8192$, and they fall into exactly four classes of $2048$:
internal ($J$-linear, $\Gamma$ kept, boosts kept), $P$ ($J$-linear, $\Gamma$ reversed, boosts reversed), $C$ ($J$-antilinear, $\Gamma$ reversed, boosts kept) and $CP$
($J$-antilinear, $\Gamma$ kept, boosts reversed). The two missing types are forbidden for every operator, since an antilinear map commuting with the Lorentz action and
keeping $\Gamma$ would identify $(\tfrac12,0)$ with its conjugate $(0,\tfrac12)$. Representatives are
\[
C=E_2E_7E_{10}E_{11}E_{12},\qquad P=O=E_1E_7E_8,\qquad CP=E_1E_2E_7E_8E_{10}E_{11}E_{12}E_{14},
\]
and all three normalise the Pati--Salam algebra and preserve the kinetic, gauge and Yukawa terms.
\end{theorem}

The factors of $C$ are instructive: $E_2$ is the surplus generator of the first constituent, where the one-constituent charge conjugation lives; $E_7$
exchanges left and right; and $E_{10}E_{11}E_{12}$, an odd number of colour generators, induces the outer automorphism of $\su4$ that exchanges $\mathbf4$ and
$\bar{\mathbf4}$. $P=O$ is parity composed with an $SU(2)$ Weyl reflection, and $CP=\pm(C\,O)\,E_7E_{14}$ with $E_7E_{14}$ a gauge element.

\begin{theorem}[$CP$ fixes the order parameter; $C$ and $P$ move it]
\label{thm:vacCP}
$CP\,\Phi_0\,CP^{\!\top}=\Phi_0$. $C\Phi_0C^{\!\top}$ is the $\nu_L\nu_L$ kernel, which commutes with $\su3$, $\su2_R$ and $Q$ but not with $\su2_L$, and
$P\Phi_0P^{\!\top}$ is the $e_Le_L$ kernel. Both images lie on the mirror orbit in $(\mathbf{10},\mathbf3,\mathbf1)$: the degeneracy of \cref{prop:LR} is realised by two distinct
discrete symmetries.
\end{theorem}

\noindent Executable check: a6-2 (exact symbolic or integer arithmetic on the $512$-dimensional module, with $\Phi_0$ rebuilt exactly; \cref{app:methods}).

Invariance of the order parameter is not invariance of the theory. The kinetic, gauge and Yukawa terms and the trace invariants are $CP$-invariant; the general
potential on $\mathcal M$ is invariant only for couplings satisfying reality conditions; and on $(\mathbf{10},\mathbf1,\mathbf3)$ the holomorphic quartic needs a generalised $CP$,
composed with a baryon-number rotation $e^{\theta(J+H/3)}$, which is a classical symmetry of the renormalisable theory for every $\kappa$ but is anomalous, since $4B$
has $SU(2)_L^2$ anomaly $2$. Whether a $CP$ survives in the quantum theory is not settled.

Which fermions are their own antiparticles follows from a pairing rule. If $M$ is a mass kernel and $t$ an unbroken generator with $[t,M]=0$, then a Majorana-type $M$
satisfies $M\hat t=-\hat tM$, so within one chirality it pairs a representation of the unbroken group with its conjugate, while a Dirac-type $M$ pairs equal
representations of opposite chirality. A single irreducible multiplet can therefore carry a Majorana mass only if its representation of the unbroken group is
real: a nonzero charge under an unbroken $U(1)$ excludes this, while a nontrivial real representation of a non-abelian group does not. In the block, every multiplet
other than $\nu_R$ carries a nonzero hypercharge at the Pati--Salam stage, and every multiplet other than $\nu_L$ and $\nu_R$ is a colour triplet or electrically
charged at the electroweak stage. So $\nu_R$ is the only Majorana fermion at the first stage, the light $\nu$ and heavy $N$ of \cref{ex:EWspectrum} are the
Majorana fermions at the second, and the massive $\nu_R$ is mapped to itself by $CP$ with eigenvalues $\pm i$, since $CP^2=-\id$ on the block.

% --- source: tex/s22_goldstone#4 (Baryon number in the potential)
\subsection{Baryon number in the potential}

For a Pati--Salam invariant $\BmL=0$, so $4B=F$ and $\Delta B=F/4$ in terms of the fermion number $F$.

\begin{theorem}[Invariant polynomials by fermion number]
\label{thm:Bcounts}
The Pati--Salam-invariant polynomials on $\mathcal M$ of degrees two and four, graded by $F$, number
\begin{center}
\small
\begin{tabular}{@{}lcccc@{}}
\toprule
degree & $F=0$ & $F=\pm4$ ($|\Delta B|=1$) & $F=\pm8$ ($|\Delta B|=2$) & total\\
\midrule
$2$ & $12$ & $3+3$, all on $2(\mathbf6,\mathbf1,\mathbf1)$ & --- & $18$\\
$4$ & $402$ & $127+127$ & $15+15$, including $\mathcal D$ & $686$\\
\bottomrule
\end{tabular}
\end{center}
With $U(1)_F$ imposed the couplings number $12+402$; with $U(1)_{1-\Gamma}$, $8+228$; with both $F$ and $J\Gamma$-charge, $6+162$.
\end{theorem}

\noindent Executable check: z4-3, a6-4, a6-5 (exact Molien--Weyl integrals; \cref{app:methods}), with the charges of the components under $J$ and $J\Gamma$
read from the realisation in a8-9 (numerical). $U(1)_{1-\Gamma}$ is generated by $J-J\Gamma$, so its invariants are those with equal charges under $J$ and
$J\Gamma$.

Pati--Salam gauge invariance alone does not forbid baryon-number violation. The six $F=\pm4$ quadratic invariants are holomorphic bilinears in the complex sextets
of $2(\mathbf6,\mathbf1,\mathbf1)$; under $\su3\oplus\uu1_{\BmL}$ each sextet is $\bar{\mathbf3}_{2/3}\oplus\mathbf3_{-2/3}$, diquark components coupling to quark pairs and leptoquark
components coupling to quark--lepton pairs, and $\epsilon^{abcd}\Delta_{ab}\Delta'_{cd}$ pairs one with the other. Through the Yukawa coupling such a mass mixing generates
tree-level $\Delta B=\Delta L=1$ four-fermion operators, of the type that mediates proton decay. Two protections are available and they are not equivalent: continuous
$U(1)_F$ forbids every $F\neq0$ term and conserves baryon number perturbatively (it is violated by $SU(2)$ instantons \cite{tHooft76}); the $\Z_4$ generated on $\mathcal M$ by
conjugation with $e^{\pi J/4}$ forbids the $|\Delta B|=1$ terms but allows the $|\Delta B|=2$ ones that mediate neutron--antineutron oscillation. $U(1)_F$, which the
radiative scenario of \cref{sec:selection} has, is sufficient, not necessary.

% --- source: tex/s22_goldstone#5 (Topology of the vacuum orbit)
\subsection{Topology of the vacuum orbit}

\begin{theorem}[$\Z_2$ strings]
\label{thm:strings}
On the block $H=3J(\BmL)$, so $e^{\theta H}$ acts on quarks as $\cos\theta+J\sin\theta$ and on leptons as $\cos3\theta-J\sin3\theta$. The colour element $g_*=e^{\pi H/3}$ fixes
$\Phi_0$, acts as $-1$ on the $\nu_R$ line and on every lepton, and the path $\theta\mapsto e^{\theta H}\Phi_0e^{-\theta H}=\cos6\theta\,\Phi_0-\sin6\theta\,J\Phi_0$,
$\theta\in[0,\pi/3]$, is a closed loop on which the condensate phase winds once. Every generator of the stabiliser annihilates the $\nu_R$ line, so the identity component
of the stabiliser acts trivially there and $g_*$ is not in it. Lifting to $SU(4)\times SU(2)_L\times SU(2)_R$, the stabiliser has exactly two components, so the gauge orbit of
the vacuum has $\pi_1=\Z_2$, and $\pi_2\supseteq\Z$ from $U(1)_Y$.
\end{theorem}

If the gauge orbit is the complete set of minima of the chosen potential, as it is on $(\mathbf{10},\mathbf1,\mathbf3)$ in the region of \cref{thm:select0}, the vacuum supports
topologically stable $\Z_2$ strings around which $\nu_R$ changes sign, the Pati--Salam form of the $\Z_2$ strings of $\so{10}$ broken by a $\mathbf{126}$ \cite{KibbleLazaridesShafi82},
and Pati--Salam monopoles. If the minimum set is larger, as in the degenerate limit with accidental $O(272)$ symmetry, the relevant fundamental group is that of the
larger manifold, and a string may unwind through flat directions. The Standard-Model orbit and its mirror lie in different Pati--Salam components, so no gauge
transformation connects them; whether they are disconnected components of the full minimum set, which would give domain walls \cite{ZeldovichKobzarevOkun75},
depends on paths through configurations outside both orbits. Along the great circle of \cref{prop:LR} the radiative energy has a barrier between them for $x<27$,
but other paths are not controlled.

